\chapter{LANDASAN TEORI}
\label{chap:teori}

Seperti telah dijelaskan pada Subbab~\ref{sec:1:latar_belakang}, lingkungan Rusunami The Jarrdin Cihampelas memiliki tingkat mobilitas penyewa yang tinggi serta variasi pola penyewaan yang cukup dinamis. Hal ini menimbulkan tantangan dalam mendeteksi potensi perilaku penyewa mencurigakan, terutama karena sistem pengawasan yang masih manual dan tidak terintegrasi di satu basis data. Dengan tersedianya data historis penyewa yang telah diberi label, pendekatan berbasis \textit{machine learning}—khususnya klasifikasi—dapat dimanfaatkan untuk mempelajari pola perilaku penyewa berisiko dan memberikan rekomendasi secara otomatis. Untuk mendukung pengembangan solusi tersebut, bab ini menyajikan kajian teori dan eksplorasi teknologi yang relevan. Pembahasan dimulai dari teori mengenai hunian vertikal dan sistem pengelolaan penyewa di lingkungan rumah susun, termasuk dinamika penyewaan dan potensi risiko yang sering muncul.

Selanjutnya, teori terkait \textit{machine learning} akan dibahas secara mendalam, meliputi paradigma pembelajaran mesin (khususnya \textit{supervised learning}), deteksi anomali (\textit{anomaly detection}), serta tahapan pemrosesan data termasuk ekstraksi fitur (\textit{feature engineering}) dan pemilihan fitur (\textit{feature selection}). Berbagai algoritma klasifikasi akan diuraikan, mencakup \textit{Decision Tree}, \textit{Random Forest}, \textit{Support Vector Machine} (SVM), \textit{K-Nearest Neighbors} (KNN), dan \textit{XGBoost}. Aspek utama dalam penelitian ini, yaitu penanganan ketidakseimbangan kelas (\textit{imbalanced dataset}), dibahas melalui teknik \textit{cost-sensitive learning} dan optimasi parameter model melalui \textit{hyperparameter tuning}. Teori mengenai metrik evaluasi model seperti \textit{precision}, \textit{recall}, \textit{F1-score}, dan \textit{confusion matrix} turut dijelaskan sebagai dasar pemilihan model terbaik untuk dataset tidak seimbang. Terakhir, bab ini juga membahas teknologi pendukung implementasi, meliputi \textbf{Python}, \textit{Streamlit}, \textbf{Node.js}, dan Google Cloud Platform (GCP). Dengan demikian, teori dan teknologi yang dipaparkan pada bab ini menjadi dasar konseptual yang digunakan dalam analisis, perancangan, implementasi, dan evaluasi sistem deteksi perilaku penyewa mencurigakan di Rusunami The Jarrdin Cihampelas.

 

% Sub-bab 2.1 rumah susun
\section{Rumah Susun}

\subsection{Definisi dan Landasan Hukum}
Rumah Susun adalah bangunan gedung bertingkat yang dibangun dalam suatu lingkungan yang terbagi dalam bagian-bagian yang distrukturkan secara fungsional, baik dalam arah horizontal maupun vertikal yang mana merupakan satuan-satuan yang masing-masing dapat dimiliki dan digunakan secara terpisah. Definisi dan regulasi utama mengenai rumah susun diatur dalam Undang-Undang Republik Indonesia Nomor 20 Tahun 2011 tentang Rumah Susun. Undang-undang ini hadir sebagai respon atas meningkatnya kebutuhan hunian di daerah perkotaan yang memiliki keterbatasan lahan (Sutedi, 2010).

Tujuan pembangunan rumah susun tidak hanya sekadar menyediakan tempat tinggal, tetapi juga untuk meningkatkan efisiensi penggunaan tanah, tata ruang, dan menciptakan lingkungan pemukiman yang lengkap, serasi, dan seimbang. Menurut Budiharjo dan Sujarto (2013) dalam bukunya \textit{Kota Berkelanjutan}, konsep hunian vertikal merupakan solusi paling rasional untuk mengatasi \textit{urban sprawl} (perembetan kota) yang tidak terkendali akibat tingginya arus urbanisasi.

\subsection{Klasifikasi Rumah Susun}
Berdasarkan peruntukan dan pengelolaannya, rumah susun dapat diklasifikasikan menjadi beberapa kategori. Menurut Undang-undang Nomor 20 Tahun 2011 tentang Rumah Susun\footnote{Undang-undang Nomor 20 Tahun 2011 tentang Rumah Susun (Lembaran Negara Republik Indonesia Tahun 2011 Nomor 108)}, klasifikasi tersebut meliputi berikut.

\begin{enumerate}
    \item \textbf{Rumah Susun Umum} \\
    Rumah susun yang diselenggarakan untuk memenuhi kebutuhan rumah bagi Masyarakat Berpenghasilan Rendah (MBR).
    \item \textbf{Rumah Susun Khusus} \\
    Rumah susun yang diselenggarakan untuk memenuhi kebutuhan khusus, seperti asrama mahasiswa atau hunian bagi lansia.
    \item \textbf{Rumah Susun Negara} \\
    Rumah susun yang dimiliki negara dan berfungsi sebagai tempat tinggal atau hunian sarana pembinaan keluarga serta penunjang pelaksanaan tugas pejabat dan/atau pegawai negeri.
    \item \textbf{Rumah Susun Komersial} \\
    Rumah susun yang diselenggarakan untuk mendapatkan keuntungan, yang sering dikenal dengan istilah {Apartemen} atau Kondominium. Objek penelitian ini, yaitu The Jarrdin Cihampelas yang mana termasuk dalam kategori Rusunami (Rumah Susun Sederhana Milik) yang dalam perkembangannya sering kali beroperasi dengan mekanisme pasar komersial.
\end{enumerate}

\subsection{Struktur Kepemilikan}
Salah satu karakteristik unik dari rumah susun yang membedakannya dengan rumah tapak (\textit{landed house}) adalah konsep kepemilikannya. Pemilik unit rumah susun tidak hanya memiliki unit huniannya secara individual, tetapi juga memiliki hak atas bagian lain dari gedung tersebut secara kolektif.

Menurut Harsono (1999) dalam bukunya berjudul \textit{Hukum Agraria Indonesia, Himpunan Peraturan-Peraturan Hukum Tanah}, kepemilikan dalam rumah susun meliputi tiga elemen tak terpisahkan sebagai berikut.
\begin{enumerate}
    \item {Hak Milik atas Satuan Rumah Susun (HMSRS)} yaitu
    hak kepemilikan yang bersifat perseorangan dan terpisah (ruang unit hunian).
    \item {Bagian bersama} yaitu 
    Bagian rumah susun yang dimiliki secara tidak terpisah untuk pemakaian bersama dalam kesatuan fungsi gedung (contoh: pondasi, kolom, dinding pemikul, atap, dan tangga).
    \item {Benda bersama} yaitu 
    Benda yang bukan merupakan bagian rumah susun melainkan bagian yang dimiliki bersama untuk pemakaian bersama (contohnya, instalasi listrik, tempat parkir, taman, dan bangunan pertokoan).
    \item {Tanah bersama} yaitu
    Sebidang tanah yang digunakan atas dasar hak bersama secara tidak terpisah yang di atasnya berdiri rumah susun.
\end{enumerate}



% \subsection{Dinamika Sosial dan Potensi Kerawanan}
% Tinggal di rumah susun membawa perubahan pola interaksi sosial dari horizontal menjadi vertikal. Marlina (2008) dalam \textit{Panduan Perancangan Bangunan Komersial} menyebutkan bahwa karakteristik penghuni apartemen atau rumah susun cenderung lebih individualis dan memiliki privasi yang tinggi dibandingkan penghuni perumahan tapak.

% Kondisi ini memiliki dua sisi mata uang. Di satu sisi, privasi penghuni sangat terjaga. Namun, di sisi lain, tingginya tingkat anonimitas (ketidaktahuan antar tetangga) dan aksesibilitas yang sering kali longgar pada apartemen sewa harian menciptakan celah keamanan. Kurangnya kontrol sosial (\textit{social control}) dari sesama penghuni memudahkan terjadinya penyalahgunaan unit untuk aktivitas ilegal, seperti prostitusi terselubung atau peredaran narkoba, karena pelaku merasa terlindungi oleh dinding privasi dan kurangnya interaksi tetangga (Soekanto, 2012). Hal inilah yang mendasari perlunya sistem pengawasan berbasis teknologi, seperti \textit{data capture} dan deteksi anomali, untuk membantu manajemen gedung dalam memitigasi risiko keamanan yang muncul akibat karakteristik sosial hunian vertikal.

% Sub-bab 2.2 pola_perilaku


% \section{Pola Perilaku Penyewa Kamar yang Berpotensi Melakukan Tindakan Terlarang}

% Tindakan terlarang dalam konteks penyewaan kamar mencakup berbagai aktivitas kriminal, khususnya penyalahgunaan narkoba dan praktik prostitusi. Penyalahgunaan narkoba mencakup penggunaan, penyimpanan, dan distribusi narkotika secara ilegal di dalam kamar sewa, termasuk berbagai jenis narkotika seperti metamfetamin, kokain, heroin, dan ekstasi sebagaimana diatur dalam Undang-Undang Nomor 35 Tahun 2009 tentang Narkotika. Sementara itu, prostitusi melibatkan eksploitasi seksual komersial yang dilakukan di kamar sewa, baik secara individu maupun dalam jaringan yang mana sering kali berkaitan dengan perdagangan manusia dan eksploitasi seksual anak (Pasal 296 dan 506 KUHP serta Undang-Undang Nomor 21 Tahun 2007 tentang Pemberantasan Tindak Pidana Perdagangan Orang).

% Regulasi hukum yang berlaku di Indonesia mencakup beberapa undang-undang utama yang berfungsi untuk mencegah aktivitas ilegal ini. Undang-Undang Nomor 35 Tahun 2009 tentang Narkotika mengatur larangan peredaran dan penggunaan narkotika secara ilegal, termasuk hukuman berat bagi pelaku. Pasal 296 dan 506 KUHP menetapkan sanksi bagi pihak yang memfasilitasi atau mengambil keuntungan dari prostitusi, sementara Undang-Undang Nomor 21 Tahun 2007 secara khusus melarang segala bentuk perdagangan manusia. Selain itu, beberapa daerah memiliki peraturan daerah (Perda) yang mengatur praktik penyewaan kamar guna mencegah tindak kriminal. 

% Di Kota Bandung, terdapat beberapa Peraturan Daerah (Perda) yang mengatur praktik penyewaan kamar dan bertujuan mencegah tindak kriminal. Salah satunya adalah Peraturan Daerah Kota Bandung Nomor 4 Tahun 2003 tentang Sewa Menyewa/Kontrak Rumah dan/atau Bangunan, yang mengatur ketentuan sewa menyewa atau kontrak rumah dan bangunan di Kota Bandung, termasuk kewajiban para pihak dan ketentuan lainnya yang bertujuan untuk menciptakan ketertiban dalam praktik penyewaan properti. Selain itu, terdapat Peraturan Daerah Kotamadya Daerah Tingkat II Bandung Nomor 07 Tahun 1986 tentang Izin Usaha Kepariwisataan yang mengatur tentang izin usaha kepariwisataan, termasuk losmen, penginapan remaja, pondok wisata, dan usaha rekreasi lainnya. Peraturan ini bertujuan untuk memastikan bahwa usaha-usaha tersebut beroperasi sesuai dengan ketentuan yang berlaku dan tidak disalahgunakan untuk aktivitas ilegal. Kemudian ada juga Peraturan Daerah Kota Bandung Nomor 9 Tahun 2019 tentang Ketertiban Umum, Ketentraman, dan Perlindungan Masyarakat yang mengatur berbagai aspek ketertiban umum di Kota Bandung, termasuk ketentuan yang dapat diterapkan untuk mencegah penyalahgunaan tempat tinggal atau kamar sewa untuk kegiatan kriminal. Dalam beberapa pasal, peraturan ini menetapkan larangan terhadap kegiatan yang mengganggu ketertiban umum yang dapat mencakup aktivitas ilegal di tempat sewa. Peraturan Daerah Kota Bandung Nomor 11 Tahun 2010 tentang Pelarangan, Pengawasan, dan Pengendalian Minuman Beralkohol mengatur larangan dan pengendalian peredaran minuman beralkohol di Kota Bandung. Meskipun fokus utamanya pada minuman beralkohol, penerapan peraturan ini dapat berkontribusi dalam mencegah tindak kriminal yang sering dikaitkan dengan konsumsi alkohol di tempat-tempat sewa. Terakhir, ada Peraturan Daerah Kota Bandung Nomor 13 Tahun 2019 tentang Perumahan, Kawasan Permukiman, dan Penanganan Kawasan Kumuh yang mengatur penyelenggaraan perumahan dan kawasan permukiman di Kota Bandung, termasuk upaya pencegahan dan peningkatan kualitas perumahan kumuh. Meskipun tidak secara spesifik mengatur penyewaan kamar, peraturan ini berkontribusi dalam menciptakan lingkungan permukiman yang tertib dan aman. Penerapan dan pengawasan terhadap Perda-Perda tersebut dilakukan oleh pemerintah daerah bersama aparat penegak hukum untuk memastikan bahwa praktik penyewaan kamar tidak disalahgunakan untuk kegiatan kriminal, termasuk di daerah Cihampelas dan wilayah lain di {Kota Bandung}.

% Penyewa yang berpotensi melakukan tindakan terlarang sering menunjukkan pola perilaku tertentu yang mencurigakan. Beberapa ciri utama termasuk penyewaan dalam jangka pendek dan tidak konsisten, sering berganti kamar atau lokasi dalam waktu yang tidak wajar, serta memiliki banyak tamu yang datang dan pergi dalam waktu singkat, terutama pada malam hari. Selain itu, pembayaran sering kali dilakukan dalam bentuk tunai dan tanpa identitas yang jelas, serta ada kecenderungan untuk menghindari interaksi dengan staf atau penghuni lain. Penyewa yang mencurigakan juga sering membawa banyak koper atau paket dalam kondisi mencurigakan, serta ditemukan alat-alat seperti jarum suntik, bong, atau bukti komunikasi digital terkait prostitusi (Leung, Yang, \& Dubin, 2018).

% Beberapa faktor yang menyebabkan hotel atau apartemen dipilih sebagai tempat tindakan terlarang antara lain kurangnya pengawasan dan keamanan, tidak adanya sistem identifikasi tamu yang ketat, serta lokasi yang strategis dengan akses mudah. Apartemen yang menerapkan sistem penyewaan fleksibel juga sering kali menjadi target utama karena tidak adanya batasan atau verifikasi bagi tamu yang sering berpindah-pindah kamar (Paraskevas \& Brookes, 2018). Selain itu, kurangnya kesadaran dan pelaporan dari pihak pengelola properti sering kali membuat aktivitas mencurigakan tidak terdeteksi.


% Untuk menangani penyewa yang melakukan tindakan kriminal, pengelola apartemen dan hotel dapat menerapkan berbagai strategi, termasuk penerapan sistem keamanan yang ketat dengan mewajibkan penyewa menunjukkan identitas asli yang diverifikasi serta menggunakan sistem CCTV dan patroli keamanan secara berkala. Pelatihan bagi staf dan manajemen sangat penting untuk mengenali pola perilaku mencurigakan serta menjalin kerja sama dengan pihak kepolisian dan organisasi terkait untuk edukasi (Zhang \& Potter, 2022). Selain itu, kerja sama dengan aparat penegak hukum perlu diperkuat dengan membentuk mekanisme pelaporan cepat dan meningkatkan koordinasi antara manajemen properti dengan kepolisian setempat. Teknologi juga dapat dimanfaatkan untuk mendeteksi aktivitas mencurigakan, misalnya melalui sistem analisis data dan {\it machine learning} yang mampu mengidentifikasi pola perilaku penyewa yang mencurigakan. Terakhir, penyuluhan dan kampanye kesadaran kepada penghuni apartemen dan staf mengenai bahaya aktivitas ilegal dapat membantu meningkatkan kewaspadaan dan mendorong pelaporan terhadap aktivitas mencurigakan.

% Hotel dan apartemen sering menjadi tempat yang dipilih untuk aktivitas ilegal seperti penyalahgunaan narkoba dan prostitusi karena berbagai faktor, termasuk lemahnya pengawasan dan regulasi penyewaan jangka pendek. Untuk mencegah dan menangani hal ini, diperlukan kombinasi strategi yang mencakup sistem keamanan yang ketat, pelatihan staf, kerja sama dengan aparat penegak hukum, serta pemanfaatan teknologi. Dengan pendekatan yang sistematis, pengelola properti dapat mengurangi risiko penyalahgunaan kamar sewa untuk aktivitas kriminal.








\section{Pola Perilaku Penyewa Kamar yang Berpotensi Melakukan Tindakan Terlarang}

Praktik penyewaan properti--seperti apartemen dan hotel-- rentan disalahgunakan untuk aktivitas ilegal, khususnya penyalahgunaan narkotika dan praktik prostitusi. Berdasarkan regulasi nasional, tindak pidana ini diatur secara ketat dan diancam dengan hukuman berat. Penyalahgunaan narkotika diatur dalam Undang-Undang Nomor 35 Tahun 2009 tentang Narkotika. Sementara itu, aktivitas yang memfasilitasi eksploitasi seksual komersial atau prostitusi dilarang keras berdasarkan Pasal 296 dan 506 KUHP, serta Undang-Undang Nomor 21 Tahun 2007 tentang Pemberantasan Tindak Pidana Perdagangan Orang (TPPO).

Di tingkat daerah, khususnya di Kota Bandung, pemerintah telah menetapkan berbagai instrumen hukum untuk memitigasi penyalahgunaan fungsi hunian dan menciptakan ketertiban umum. Beberapa regulasi utama tersebut sebagai berikut.
\begin{itemize}
    \item \textbf{Perda Kota Bandung No. 4 Tahun 2003} tentang Sewa Menyewa/Kontrak Rumah dan/atau Bangunan, yang mengatur kewajiban para pihak untuk menjaga ketertiban praktik penyewaan properti.
    \item \textbf{Perda Kotamadya Dati II Bandung No. 07 Tahun 1986} tentang Izin Usaha Kepariwisataan, guna memastikan penginapan dan losmen beroperasi sesuai izin dan tidak disalahgunakan.
    \item \textbf{Perda Kota Bandung No. 11 Tahun 2010} tentang Pelarangan, Pengawasan, dan Pengendalian Minuman Beralkohol, yang membatasi pemicu tindak kriminal di tempat sewa.
    \item \textbf{Perda Kota Bandung No. 9 Tahun 2019} (Ketertiban Umum) dan \textbf{Perda No. 13 Tahun 2019} (Perumahan dan Kawasan Permukiman), yang memberikan dasar hukum penindakan terhadap aktivitas yang mengganggu ketentraman masyarakat di area permukiman.
\end{itemize}

Menurut literatur kriminologi dan perhotelan, penyewa yang berpotensi melakukan tindakan terlarang umumnya menunjukkan pola perilaku (\textit{behavioral patterns}) spasial dan temporal yang anomali. Leung, Yang, dan Dubin (2018) mengidentifikasi beberapa indikator perilaku mencurigakan, antara lain penyewaan dalam jangka waktu sangat pendek (\textit{micro-stays}) yang sering dilakukan berulang; intensitas keluar-masuk tamu yang tinggi dalam waktu singkat, terutama pada malam hari; kecenderungan membayar menggunakan uang tunai untuk menghindari pelacakan identitas; serta keengganan berinteraksi dengan staf atau penghuni lain (\textit{avoidance behavior}). Selain itu, kecurigaan fisik sering kali didukung dengan penemuan barang bawaan tak lazim atau residu aktivitas ilegal (seperti jarum suntik atau bong).

Tingginya kerentanan apartemen terhadap aktivitas ini disebabkan oleh karakteristik struktural penyewaan itu sendiri. Paraskevas dan Brookes (2018) dalam kajian keamanan perhotelan menyatakan bahwa properti yang menerapkan sistem sewa fleksibel menjadi target utama karena lemahnya verifikasi identitas, minimnya pengawasan langsung (\textit{blind spots}), serta tingkat privasi yang tinggi yang menyulitkan deteksi dini oleh pengelola. 

Untuk menangani ancaman tersebut, pengelola properti memerlukan pendekatan keamanan yang berlapis (\textit{multilayered security}). Pendekatan ini mencakup penerapan verifikasi identitas otentik, optimalisasi pengawasan CCTV, dan pelatihan staf untuk mengenali anomali perilaku tamu (Zhang \& Potter, 2022). Lebih lanjut, adopsi teknologi berbasis data dan \textit{machine learning} menjadi esensial di era modern. Sistem deteksi otomatis dapat menganalisis rekam jejak transaksi, durasi inap, dan pola kedatangan penyewa untuk menghasilkan peringatan dini (\textit{early warning system}) bagi pihak manajemen. Melalui kombinasi antara protokol keamanan fisik, kolaborasi dengan aparat penegak hukum, dan sistem algoritma cerdas, risiko penyalahgunaan kamar sewa untuk aktivitas kriminal dapat direduksi secara signifikan.

% Sub-bab 2.3 machine_learning
\section{\textit{Machine Learning}}\label{sec:machine_learning}

Dalam bidang \textit{machine learning}, klasifikasi merupakan salah satu tugas utama yang bertujuan untuk memetakan input data ke dalam label kelas tertentu berdasarkan pola yang dipelajari dari data. Secara umum, terdapat empat paradigma utama dalam membangun model pembelajaran mesin, yaitu \textit{unsupervised learning}, \textit{semi-supervised learning}, \textit{supervised learning}, dan \textit{reinforcement learning} (Murphy, 2012; Kelleher et al., 2015; Sutton \& Barto, 2018).

\textit{Unsupervised learning} digunakan ketika data tidak memiliki label kelas, sehingga algoritma berusaha menemukan struktur, pola, atau distribusi tersembunyi dalam data secara mandiri. Sebagai contoh, algoritma \textit{clustering} dapat mengelompokkan pelanggan toko berdasarkan kemiripan kebiasaan belanja mereka tanpa perlu mengetahui kategori pelanggan tersebut sebelumnya. Contoh lain adalah mendeteksi anomali menggunakan algoritma seperti \textit{Isolation Forest}, \textit{Local Outlier Factor}, dan \textit{One-Class SVM} (Chandola, Banerjee, \& Kumar, 2009; Aggarwal, 2013). Teknik ini umum digunakan dalam kasus ketika tidak tersedia informasi eksplisit mengenai kelas data, misalnya untuk memisahkan profil penyewa yang normal dan yang mencurigakan tanpa adanya data riwayat pelanggaran sebelumnya.

\textit{Semi-supervised learning} digunakan pada dataset yang memiliki sedikit data berlabel dan mayoritas data tidak berlabel. Metode ini memanfaatkan sejumlah kecil data berlabel tersebut untuk memandu proses pembelajaran dan mengekstrapolasi label ke data yang belum dilabeli. Pendekatan ini sering digunakan ketika proses pelabelan manual memerlukan waktu dan biaya pakar yang sangat tinggi (Chapelle et al., 2006). Contohnya adalah pengkategorian artikel berita, di mana hanya sebagian kecil artikel yang diberikan label topik (misalnya ``olahraga'' atau ``politik'') oleh manusia, dan algoritma menggunakan pola kata dari artikel tersebut untuk melabeli jutaan artikel lainnya secara otomatis (Zhu, 2005).

Sementara itu, \textit{supervised learning} digunakan ketika seluruh data latih telah dilabeli dengan kelas target tertentu. Model dilatih menggunakan pasangan \textit{input-output} (fitur dan label), sehingga mampu mempelajari pemetaan matematis antar fitur dan menghasilkan prediksi pada data baru (James et al., 2013). Contohnya adalah sistem penyaring email \textit{spam}, di mana model dilatih menggunakan ribuan riwayat email yang sudah ditandai secara eksplisit sebagai ``spam'' atau ``bukan spam'' agar mampu mengklasifikasikan email yang baru masuk. Algoritma yang umum digunakan dalam pendekatan ini meliputi \textit{Decision Tree, Random Forest, Support Vector Machine, K-Nearest Neighbors,} dan \textit{XGBoost} (Tan, Steinbach, \& Kumar, 2019).

Di sisi lain, \textit{reinforcement learning} merupakan pendekatan prediktif, di mana sebuah agen (program) belajar mengambil serangkaian keputusan melalui interaksi langsung dengan lingkungannya untuk memaksimalkan sinyal imbalan (\textit{reward}) (Sutton \& Barto, 2018). Algoritma ini tidak diajari tindakan mana yang benar melalui label data secara statis, melainkan harus mengeksplorasi dan menemukan sendiri tindakan yang menghasilkan \textit{reward} tertinggi melalui mekanisme coba-coba (\textit{trial-and-error}). Contoh penerapannya adalah pada kecerdasan buatan yang belajar bermain catur; agen tersebut akan mendapat imbalan (+1) saat memakan bidak lawan, dan mendapat hukuman atau penalti (-1) saat membuat langkah yang merugikan (Silver et al., 2016).

Dalam proses pelatihan model \textit{machine learning}, salah satu tantangan terbesar yang sering terjadi adalah \textit{overfitting}. \textit{Overfitting} merupakan kondisi di mana model menghafal data pelatihan dengan terlalu detail, termasuk menangkap gangguan kecil (\textit{noise}) dan variasi acak yang tidak penting (Dietterich, 1995). Akibatnya, model tampak memiliki tingkat akurasi yang sangat luar biasa tinggi pada data latih, namun ketika diuji menggunakan data baru yang belum pernah dilihat sebelumnya, kemampuannya menurun secara drastis (Hawkins, 2004). Analogi sederhananya seperti seorang siswa yang hanya menghafal soal-soal dan jawaban dari buku latihan ujian tanpa benar-benar memahami konsep dasar materinya. Saat menghadapi ujian sebenarnya dengan angka yang sedikit diubah, ia akan gagal menjawab (Goodfellow, Bengio, \& Courville, 2016). Untuk mengatasi masalah \textit{overfitting}, beberapa strategi umum yang diterapkan meliputi penggunaan lebih banyak data pelatihan, membatasi kerumitan model, ataupun menggunakan metode regularisasi (Hastie, Tibshirani, \& Friedman, 2009).

% Berdasarkan karakteristik data yang telah dilabeli sebagai normal atau tidak normal, maka pendekatan yang digunakan dalam penelitian ini adalah \textit{supervised learning}. Hal ini memungkinkan pelatihan model klasifikasi yang dapat mengenali pola perilaku berdasarkan data historis, serta mengevaluasi performanya secara objektif menggunakan metrik evaluasi seperti akurasi, presisi, dan recall.





% Sub-bab 2.4 Anomaly Detection
% \newpage
\section{Deteksi Anomali}

{Deteksi anomali} merupakan proses identifikasi pola atau data yang menyimpang secara signifikan dari perilaku mayoritas dan dapat digunakan untuk mengungkap aktivitas tidak biasa seperti penipuan, intrusi sistem, maupun penyewa mencurigakan (Chandola, Banerjee, \& Kumar, 2009; Aggarwal, 2013). Deteksi anomali berguna ketika data yang dianalisis tidak memiliki label, sehingga membutuhkan metode {\it unsupervised learning}. Dua metode umum dalam kategori ini adalah {\it Isolation Forest} dan {\it One-Class Support Vector Machine (SVM)}.
    
    \begin{enumerate}
      \item {{\it Isolation Forest}}
    
      {\it Isolation Forest} adalah algoritma berbasis pohon keputusan yang didesain untuk mengisolasi data anomali lebih cepat daripada data normal dengan cara membuat pembelahan acak berulang terhadap fitur data (Liu, Ting, \& Zhou, 2008). Data yang dapat dipisahkan dalam sedikit pembelahan dianggap sebagai anomali. Kedalaman rata-rata isolasi sebuah data dari banyak pohon digunakan untuk menghitung skor anomali. Berikut rumusnya
      \[
      s(x, n) = 2 - \frac{E(h(x))}{c(n)}
      \]

      \[
        s(x, n) = 2^{-\frac{E(h(x))}{c(n)}}
      \]
        
      dengan \(E(h(x))\) adalah rata-rata kedalaman isolasi dari data \(x\), dan \(c(n) \approx \ln(n) + 0{.}5772\) adalah nilai koreksi kedalaman untuk dataset berukuran \(n\) (nilai eksak untuk harmonik ke-\(n\)).
    
      Misalkan data penyewa A memiliki kedalaman isolasi 2, 3, dan 2 dari 3 pohon yang dibentuk, maka:
      \[
      E(h(x)) = \frac{2 + 3 + 2}{3} = 2{,}33
      \quad \text{dan} \quad c(n) = 1{.}5 \ (\text{asumsi dataset kecil})
      \]
      Maka skor anomali:
      \[
        s(x, n) = 2^{-\frac{2,33}{1,5}} = 2^{-1,553} \approx 0,341
      \]
      Skor mendekati 0{.}5 menunjukkan data cenderung normal, sedangkan jika skor mendekati 1, maka data dianggap anomali kuat (Liu et al., 2008).
    
      \item {{\it One-Class SVM}}
    
        {\it One-Class SVM} adalah metode margin maksimal yang dirancang untuk memisahkan distribusi data normal dari titik asal ({\it origin}) dalam ruang fitur (Sch\"{o}lkopf et al., 2001). Algoritma ini sering menggunakan fungsi kernel seperti {\it Radial Basis Function} (RBF) untuk memetakan data ke dimensi yang lebih tinggi, berikut rumusnya
        \[
        K(x_i, x_j) = \exp\left(-\gamma \cdot \|x_i - x_j\|^2\right)
        \]
        dengan \(\gamma\) sebagai parameter lebar kernel dan \(\|x_i - x_j\|^2\) sebagai kuadrat jarak Euclidean.
        
        Misalkan dua data penyewa \(A = (1, 2)\) dan \(B = (2, 3)\) dengan \(\gamma = 0{,}5\), maka:
        \[
        \|A - B\|^2 = (1-2)^2 + (2-3)^2 = 2
        \quad \Rightarrow \quad
        K(A, B) = \exp(-0{,}5 \times 2) \approx 0{,}368
        \]
        
        Nilai kernel ini merepresentasikan skor kemiripan antar data. Jika skor mendekati 1, berarti kedua data tersebut sangat mirip. Sebaliknya, jika skor kecil atau mendekati 0 (seperti hasil \(0{,}368\) pada contoh), maka data tersebut jauh berbeda. Dalam konteks deteksi anomali, jika ada penyewa yang memiliki skor kemiripan sangat rendah terhadap mayoritas penyewa lain, maka akan ditandai sebagai pencilan ({\it outlier}).
        \end{enumerate}

% Sub-bab 2.5 feature_engineering
\section{Feature Engineering}\label{sec:feature_engineering}

\textit{Feature engineering} adalah proses membuat fitur-fitur (kolom data) baru atau mengubah fitur yang sudah ada agar representasi data menjadi lebih informatif bagi algoritma \textit{machine learning}. Menurut Kuhn dan Johnson (2019) dalam \textit{Feature Engineering and Selection: A Practical Approach for Predictive Models}, \textit{feature engineering} yang efektif sering kali menjadi faktor penentu utama keberhasilan sebuah model prediktif, dikarenakan algoritma tidak dapat menemukan pola yang kompleks jika data mentahnya tidak merepresentasikan masalah dengan baik. Tujuannya adalah membantu model untuk belajar dari data dan mengidentifikasi pola yang menyimpang dari kondisi normal (anomali).

Berdasarkan data operasional sistem penyewaan, \textit{feature engineering} berfokus pada ekstraksi fitur temporal (waktu) dan pemrosesan informasi kualitatif (teks) guna mendeteksi pola perilaku penyewa yang mencurigakan. Dalam konteks data penyewaan, berikut adalah contoh perhitungan \textit{feature engineering} dari fitur-fitur utama yang dikembangkan beserta landasan teorinya.

\begin{enumerate}
    \item \textbf{Indikator Pelanggaran Berdasarkan Komentar (\textit{Is Melanggar})} \\
    Dalam sistem pencatatan operasional, anomali sering kali tidak tercatat dalam bentuk metrik angka, melainkan melalui catatan bebas (\textit{free-text}) atau log laporan dari staf \cite{aggarwal2017outlier}. Jika kolom komentar pada riwayat penyewa tidak kosong (\textit{not null}), sistem akan menghasilkan fitur biner baru bernilai 1 (\textit{True}), dan 0 (\textit{False}) jika sebaliknya. Hal ini berfungsi sebagai sinyal kuat bagi algoritma mengenai riwayat pelanggaran tata tertib.

    \item \textbf{Durasi Sewa dalam Jam (\textit{Duration Hours})} \\
    Fitur ini menghitung total durasi penyewaan secara kontinu dalam satuan jam, berdasarkan selisih antara waktu \textit{check-out} dan waktu \textit{check-in}.
    
    {Contoh kasus perhitungan:}
    \begin{center}
        \begin{tabular}{|c|c|c|}
            \hline
            \textbf{Waktu Check-in} & \textbf{Waktu Check-out} \\
            \hline
            2025-05-18 20:00:00 & 2025-05-18 23:00:00\\
            \hline
            2025-05-19 14:00:00 & 2025-05-20 12:00:00\\
            \hline
        \end{tabular}
    \end{center}
    
    Persamaan umum yang digunakan adalah:
    \[
    \text{Duration Hours} = \frac{(\text{Check-out} - \text{Check-in})_{\text{detik}}}{3600}
    \]
    
    \begin{itemize}
        \item \textbf{Kasus 1 (Penyewaan di hari yang sama):} \\
        Selisih waktu (20:00:00 hingga 23:00:00) adalah 3 jam. Dalam satuan detik bernilai $3 \times 3600 = 10.800$ detik.
        \[
        \text{Durasi}_1 = \frac{10.800}{3600} = 3{,}0 \text{ jam}
        \]
        \item \textbf{Kasus 2 (Penyewaan melintasi hari):} \\
        Selisih waktu (19 Mei 14:00:00 hingga 20 Mei 12:00:00) adalah 22 jam. Dalam satuan detik bernilai $22 \times 3600 = 79.200$ detik.
        \[
        \text{Durasi}_2 = \frac{79.200}{3600} = 22{,}0 \text{ jam}
        \]
    \end{itemize}

    Fitur ini krusial untuk mendeteksi penyewa ``transit''. Baesens, Van Vlasselaer, dan Verbeke  \cite{baesens2015fraud} menyatakan bahwa ekstraksi fitur temporal adalah teknik paling efektif dalam sistem deteksi anomali. Sejalan dengan temuan tersebut, Gupta et al.  \cite{gupta2014outlier} menegaskan bahwa anomali temporal sering muncul dalam bentuk durasi kejadian yang terlampau singkat. Dalam konteks sewa properti, durasi inap yang sangat pendek (\textit{micro-stays}) berkorelasi positif dengan aktivitas ilegal.

    \item \textbf{Ekstraksi Waktu Check-in (\textit{Hour} \& \textit{Day of Week})} \\
    Fitur ini memecah informasi waktu kedatangan menjadi komponen jam operasional (0--23) dan indeks hari dalam seminggu (0 = Senin, 6 = Minggu). Perilaku manusia memiliki pola musiman dan temporal yang dapat diprediksi \cite{zheng2014urban}. Kedatangan pada jam tidak lazim (misalnya pukul 03.00 dini hari) dapat menjadi indikator awal (\textit{early warning}) dari profil anomali jika digabungkan dengan fitur durasi yang singkat.

    \item \textbf{Indikator Akhir Pekan (\textit{Is Weekend Check-in})} \\
    Fitur biner ini menandai apakah penyewaan dilakukan pada akhir pekan (Sabtu atau Minggu). Menurut Aggarwal  \cite{aggarwal2017outlier}, waktu merupakan atribut kontekstual (\textit{contextual attribute}) yang mendefinisikan batasan normalitas suatu data. Lonjakan penyewaan di akhir pekan adalah hal wajar, namun bisa menjadi sinyal mencurigakan jika terjadi di pertengahan minggu kerja. Chandola, Banerjee, dan Kumar  \cite{chandola2009anomaly} mendefinisikan fenomena ini sebagai anomali kontekstual. Fitur ini membantu model menghindari kesalahan klasifikasi (\textit{false positive}) dengan memberikan konteks waktu yang tepat.
\end{enumerate}

Selain melakukan ekstraksi fitur seperti waktu dan hari, tahap persiapan data untuk \textit{machine learning} harus mempertimbangkan rentang skala matematis dari data asli. Algoritma seperti \textit{Support Vector Machine} (SVM) sangat sensitif terhadap nilai fitur yang memiliki variasi skala berbeda-beda karena mengandalkan kalkulasi jarak numerik antar titik data. Oleh karena itu, fitur numerik harus distandardisasi. Menurut Bishop  \cite{bishop2006pattern} dalam \textit{Pattern Recognition and Machine Learning}, standardisasi (\textit{StandardScaler}) beroperasi dengan menyesuaikan setiap nilai fitur hingga memiliki rata-rata nol ($\mu = 0$) dan standar deviasi satu ($\sigma = 1$). Persamaan matematika standardisasi adalah $z = \frac{x - \mu}{\sigma}$, di mana fitur hasil standardisasi dipusatkan agar algoritma konvergen lebih cepat dan secara geometris proporsional tanpa ada satu fitur bernilai ribuan yang mendominasi fitur bernilai desimal kecil.

Selama mengeksekusi tahapan standardisasi maupun manipulasi fitur lainnya, kehati-hatian secara ilmiah harus dikedepankan guna mencegah fenomena kebocoran data (\textit{data leakage}). Kaufman et al.  \cite{kaufman2011leakage} mendefinisikan \textit{data leakage} sebagai suatu kelalaian di mana informasi krusial dari kumpulan data pengujian (\textit{testing set}) diam-diam masuk, membaur, atau ikut terhitung saat proses pembentukan model pada data pelatihan (\textit{training set}). Apabila misalnya nilai rata-rata keseluruhan (gabungan rasio data latih dan pengujian) dipakai pada proses \textit{StandardScaler}, model akan mengingat pola dari data masa depan, sehingga membuat akurasinya tampak fantastis secara keliru, namun akan gagal secara drastis saat membedakan kasus anomali baru karena model gagal belajar menggeneralisasi. 

% Berikut adalah contoh implementasi kode Python menggunakan pustaka \texttt{pandas} untuk mengeksekusi proses \textit{feature engineering} tersebut.

% \begin{lstlisting}[language=Python, caption={Implementasi Feature Engineering Menggunakan Python}]
% # Pastikan kolom yang dibutuhkan tersedia di dataframe
% cols_to_check = [
%     'tanggal_checkin', 'waktu_checkin', 'tanggal_checkout', 'waktu_checkout',
%     'lama_menginap', 'lantai', 'tower', 'status_kewarganegaraan', 
%     'metode_pembayaran', 'jenis_sewa', 'komentar' 
% ]

% for col in cols_to_check:
%     if col not in df.columns:
%         df[col] = np.nan
        
% # Fitur Indikator Pelanggaran (Berdasarkan Komentar)
% # Menghasilkan nilai 1 (True) jika komentar ada dan bukan string kosong
% df['melanggar'] = df['komentar'].notna() & (df['komentar'].str.strip() != '')
% df['melanggar'] = df['melanggar'].astype(int)

% # Fungsi Penggabungan Tanggal dan Waktu
% def combine_datetime(date_col, time_col):
%     dt_str = df[date_col].astype(str).str.strip() + ' ' + df[time_col].astype(str).str.strip()
%     return pd.to_datetime(dt_str, errors='coerce')

% df['checkin_dt'] = combine_datetime('tanggal_checkin', 'waktu_checkin')
% df['checkout_dt'] = combine_datetime('tanggal_checkout', 'waktu_checkout')

% # Perhitungan Durasi Sewa dalam Jam (Duration Hours)
% df['duration_hours'] = (df['checkout_dt'] - df['checkin_dt']).dt.total_seconds() / 3600
% # Mengisi nilai kosong dengan 0 dan mencegah nilai negatif (kesalahan input)
% df['duration_hours'] = df['duration_hours'].fillna(0).clip(lower=0)

% # Ekstraksi Fitur Waktu (Jam dan Hari)
% df['checkin_hour'] = df['checkin_dt'].dt.hour.fillna(0)       # Jam (0-23)
% df['checkout_hour'] = df['checkout_dt'].dt.hour.fillna(0)     # Jam (0-23)
% df['checkin_dow'] = df['checkin_dt'].dt.dayofweek.fillna(-1)  # Hari (0=Senin, 6=Minggu)

% # Fitur Biner Akhir Pekan
% df['is_weekend_checkin'] = df['checkin_dow'].isin([5, 6]).astype(int)

% # Normalisasi Tipe Data
% df['lama_menginap'] = pd.to_numeric(df['lama_menginap'], errors='coerce').fillna(0)
% \end{lstlisting}

% \textbf{Penjelasan Kode Program:}
% \begin{itemize}
%     \item \texttt{df['komentar'].notna()}: Fungsi untuk mendeteksi baris data yang memiliki entri kualitatif (bukan \textit{null} atau \textit{NaN}). Hasil logika \textit{boolean} kemudian dikonversi menjadi numerik melalui \texttt{astype(int)} agar kompatibel dengan model \textit{machine learning}.
%     \item \texttt{combine\_datetime()}: Menggabungkan variabel tanggal dan waktu yang awalnya terpisah (bertipe \textit{string}) menjadi satu objek tipe \texttt{datetime} terpadu untuk operasi matematika waktu.
%     \item \texttt{.dt.total\_seconds() / 3600}: Mengonversi nilai \textit{timedelta} (selisih waktu absolut) menjadi representasi numerik desimal dalam satuan jam (\textit{hours}).
%     \item \texttt{.clip(lower=0)}: Sebuah fungsi proteksi untuk membatasi batas bawah matriks menjadi nol, guna mencegah kemunculan durasi negatif akibat anomali kesalahan ketik saat staf memasukkan data \textit{check-in} dan \textit{check-out}.
% \end{itemize}


% Sub-bab 2.6 feature_selection
\section{Feature Selection}\label{sec:feature_selection}

Proses seleksi fitur (\textit{feature selection}) bertujuan untuk memilih bagian data (kolom) yang paling berkaitan dengan target akhir. Pemilihan ini penting untuk membuat hasil prediksi lebih akurat, mempercepat proses perhitungan komputer, dan mencegah model hanya menghafal data tanpa memahami polanya (\textit{overfitting}). Secara umum, pemilihan fitur dibagi menjadi tiga cara, yaitu berdasarkan pengetahuan dari pakar (\textit{domain knowledge}), perhitungan statistik (\textit{filter methods}), dan pendekatan heuristik atau pencarian terarah (\textit{wrapper methods}).

\subsection{Pendekatan Pengetahuan Pakar (\textit{Domain Knowledge})}
Sebelum melakukan pemilihan data kepada algoritma komputer, langkah pertama yang paling penting adalah menyaring data menggunakan pengetahuan pakar di lapangan (\textit{domain knowledge}). Cara ini dilakukan dengan menentukan informasi apa saja yang secara pengalaman nyata berkaitan erat dengan apa yang ingin ditebak, misalnya melalui wawancara mendalam dengan Pengelola apartemen atau petugas keamanan.

Menurut Kuhn dan Johnson  \cite{kuhn2013applied}, pemanfaatan pengetahuan domain sering kali lebih penting daripada teknik pemilihan fitur otomatis mana pun. Hal ini dikarenakan algoritma seleksi fitur murni secara matematis berpotensi memilih variabel yang hanya memiliki korelasi kebetulan (\textit{spurious correlation}) tanpa adanya dasar hubungan sebab-akibat yang logis. Untuk mendeteksi penyewa yang bermasalah, wawasan dari para pakar lapangan ini sangat membantu menentukan fitur khusus---seperti ada tidaknya riwayat pelanggaran atau ketidakwajaran waktu \textit{check-in}---sebagai tanda yang kuat (\textit{strong signal}) adanya aktivitas yang tidak normal.

\subsection{Metode Statistik (Filter Methods)}
Setelah fitur-fitur awal ditentukan, metode statistik digunakan untuk menghitung seberapa kuat hubungan antara fitur tersebut dengan status penyewa secara matematis. Mengingat pada kasus ini jumlah penyewa yang tidak normal (1) jauh lebih sedikit dibandingkan yang normal (0), perhitungan berikut merupakan adaptasi konseptual dari buku dasar \textit{Data Mining} \cite{han2011data} dan teori informasi \cite{cover2006elements} yang disimulasikan menggunakan perbandingan data tidak seimbang (\textit{imbalanced}).
    
\begin{enumerate}
    \item \textbf{Chi-Square Test} \\
    Uji Chi-Square ($\chi^2$) digunakan untuk melihat seberapa kuat hubungan antara dua data berbentuk kategori, seperti metode pembayaran dengan status penyewa. Menurut \cite{han2011data}, rumusnya adalah sebagai berikut:
    $$
    \chi^2 = \sum \frac{(O_i - E_i)^2}{E_i}
    $$
    dengan:
    \begin{itemize}
        \item $O_i$ : nilai observasi (angka yang sebenarnya terjadi),
        \item $E_i$ : nilai harapan (angka teoretis jika tidak ada hubungan).
    \end{itemize}

    Misalkan kita ingin menguji hubungan metode pembayaran (tunai/kartu) dengan status penyewa (normal/tidak normal) di mana kelas tidak normal merupakan minoritas.

    \begin{table}[H]
    \centering
    \caption{Tabel Kontingensi Metode Pembayaran dan Status Penyewa}
        \begin{tabular}{|c|c|c|c|}
            \hline
            \textbf{Metode Pembayaran} & \textbf{Normal (0)} & \textbf{Tidak Normal (1)} & \textbf{Total} \\
            \hline
            Tunai & 40 & 15 & 55 \\
            Kartu & 40 & 5 & 45 \\
            \hline
            Total & 80 & 20 & 100 \\
            \hline
        \end{tabular}
    \end{table}

    Langkah-langkah perhitungannya adalah sebagai berikut. Nilai observasi ($O_{ij}$) seperti angka 40, 15, dan 5 didapatkan langsung dari perhitungan data riil pada tabel kontingensi di atas. Kemudian, nilai harapan ($E_{ij}$) untuk masing-masing sel dihitung menggunakan rumus:
    $$
    E_{ij} = \frac{\text{(Total Baris)} \times \text{(Total Kolom)}}{\text{Grand Total}}
    $$
    
    Berikut adalah rincian perhitungan nilai harapan ($E_i$) untuk setiap kombinasi kategori:
    \begin{itemize}
        \item \textbf{Tunai dan Normal}:
        $$
        E_{\text{Tunai, Normal}} = \frac{55 \times 80}{100} = 44
        $$
        \item \textbf{Tunai dan Tidak Normal}:
        $$
        E_{\text{Tunai, Tidak Normal}} = \frac{55 \times 20}{100} = 11
        $$
        \item \textbf{Kartu dan Normal}:
        $$
        E_{\text{Kartu, Normal}} = \frac{45 \times 80}{100} = 36
        $$
        \item \textbf{Kartu dan Tidak Normal}:
        $$
        E_{\text{Kartu, Tidak Normal}} = \frac{45 \times 20}{100} = 9
        $$
    \end{itemize}

    Berdasarkan rincian nilai $O_i$ dan $E_i$ tersebut, perhitungan total nilai $\chi^2$ adalah sebagai berikut.
    $$
    \chi^2 = \frac{(40 - 44)^2}{44} + \frac{(15 - 11)^2}{11} + \frac{(40 - 36)^2}{36} + \frac{(5 - 9)^2}{9}
    $$
    $$
    \chi^2 \approx 0{,}36 + 1{,}45 + 0{,}44 + 1{,}77 = 4{,}02
    $$

    Semakin besar angka $\chi^2$, semakin kuat hubungan antara fitur tersebut dengan target hasil. Nilai $p$ dari distribusi hasil ini kemudian bisa digunakan untuk memastikan apakah hubungan tersebut benar-benar dapat dipercaya secara statistik (signifikan).

    \item \textbf{Korelasi Pearson} \\
    Korelasi Pearson digunakan untuk mengukur hubungan berupa garis lurus (linier) antara data berbentuk angka (misalnya lama menginap) dengan status penyewa \cite{han2011data}. Berikut rumusnya.
    $$
    r = \frac{\sum (x_i - \bar{x})(y_i - \bar{y})}{\sqrt{\sum (x_i - \bar{x})^2 \sum (y_i - \bar{y})^2}}, \quad r \in [-1,1]
    $$
    
    Sebagai contoh, pada sampel acak berikut penyewa tidak normal (1) muncul lebih sedikit:
    \begin{table}[H]
    \centering
    \caption{Data Lama Menginap dan Status}
    \begin{tabular}{|c|c|}
    \hline
    \textbf{Lama Menginap} ($x$) & \textbf{Status Tidak Normal} ($y$) \\
    \hline
    2 & 0 \\
    3 & 0 \\
    5 & 0 \\
    1 & 1 \\
    10 & 0 \\
    \hline
    \end{tabular}
    \end{table}
    
    \begin{enumerate}
        \item Hitung rata-rata untuk $x$ dan $y$.
        $$
        \bar{x} = \frac{2 + 3 + 5 + 1 + 10}{5} = \frac{21}{5} = 4{,}2, \quad
        \bar{y} = \frac{0 + 0 + 0 + 1 + 0}{5} = \frac{1}{5} = 0{,}2
        $$
    
        \item Hitung selisih nilai dari rata-rata, kuadratnya, dan hasil kali.
        \begin{table}[H]
        \centering
        \caption{Perhitungan Komponen Korelasi Pearson}
        \begin{tabular}{|c|c|c|c|c|c|c|}
        \hline
        $x_i$ & $y_i$ & $x_i - \bar{x}$ & $y_i - \bar{y}$ & $(x_i - \bar{x})^2$ & $(y_i - \bar{y})^2$ & $(x_i - \bar{x})(y_i - \bar{y})$ \\
        \hline
        2 & 0 & -2{,}2 & -0{,}2 & 4{,}84 & 0{,}04 & 0{,}44 \\
        3 & 0 & -1{,}2 & -0{,}2 & 1{,}44 & 0{,}04 & 0{,}24 \\
        5 & 0 & 0{,}8 & -0{,}2 & 0{,}64 & 0{,}04 & -0{,}16 \\
        1 & 1 & -3{,}2 & 0{,}8 & 10{,}24 & 0{,}64 & -2{,}56 \\
        10 & 0 & 5{,}8 & -0{,}2 & 33{,}64 & 0{,}04 & -1{,}16 \\
        \hline
        \multicolumn{4}{|c|}{\textbf{Total}} & 50{,}80 & 0{,}80 & -3{,}20 \\
        \hline
        \end{tabular}
        \end{table}
    
        \item Hitung nilai Pearson $r$.
        $$
        r = \frac{-3{,}20}{\sqrt{50{,}80 \times 0{,}80}} = \frac{-3{,}20}{\sqrt{40{,}64}} \approx \frac{-3{,}20}{6{,}37} \approx -0{,}502
        $$
    \end{enumerate}
    
    Nilai $r = -0{,}502$ menunjukkan adanya hubungan negatif berskala sedang antara lama menginap dan status pada kelompok data ini. Artinya, semakin singkat seseorang menginap, probabilitas ia dicurigai sebagai penyewa tidak normal cenderung semakin tinggi.
    
    \item \textbf{Mutual Information (MI)} \\
    Berdasarkan teori informasi \cite{cover2006elements}, MI mengukur seberapa banyak informasi (petunjuk) yang bisa didapatkan tentang target akhir jika kita mengetahui sebuah fitur tertentu, baik hubungannya berupa garis lurus maupun tidak. Metode ini cocok untuk data angka yang sudah dikelompokkan. Berikut rumusnya:
    $$
    I(X;Y) = \sum_{x \in X} \sum_{y \in Y} p(x,y) \cdot \log_2 \left( \frac{p(x,y)}{p(x)p(y)} \right)
    $$
    dengan:
    \begin{itemize}
        \item $p(x,y)$ : peluang kemunculan bersama antara fitur dan target,
        \item $p(x), p(y)$ : peluang kemunculan masing-masing secara terpisah.
    \end{itemize}
    
    Misalkan fitur lama menginap dikelompokkan menjadi: $\le$ 2 hari, 3--5 hari, dan $>$ 5 hari. Distribusi populasinya didominasi oleh kelas normal (0):
    
    \begin{table}[H]
    \centering
    \caption{Frekuensi Gabungan Lama Menginap dan Status Penyewa}
    \begin{tabular}{|c|c|c|}
    \hline
    \textbf{Lama menginap} & \textbf{Status} & \textbf{Frekuensi} \\
    \hline
    $\leq$ 2 hari & Normal (0) & 25 \\
    $\leq$ 2 hari & Tidak Normal (1) & 15 \\
    3--5 hari & Normal (0) & 35 \\
    3--5 hari & Tidak Normal (1) & 5 \\
    $>$ 5 hari & Normal (0) & 15 \\
    $>$ 5 hari & Tidak Normal (1) & 5 \\
    \hline
    \multicolumn{2}{|c|}{\textbf{Total}} & 100 \\
    \hline
    \end{tabular}
    \end{table}
    
    \begin{enumerate}
        \item Hitung probabilitas gabungan $p(x, y)$:
        $$
        \begin{aligned}
        p(\leq2, \text{Normal}) &= 0{,}25, \quad &p(\leq2, \text{Tidak normal}) &= 0{,}15 \\
        p(3{-}5, \text{Normal}) &= 0{,}35, \quad &p(3{-}5, \text{Tidak normal}) &= 0{,}05 \\
        p(>5, \text{Normal}) &= 0{,}15, \quad &p(>5, \text{Tidak normal}) &= 0{,}05
        \end{aligned}
        $$
    
        \item Hitung probabilitas individual:
        $$
        \begin{aligned}
        p(\leq2) &= 0{,}40, \quad p(3{-}5) = 0{,}40, \quad p(>5) = 0{,}20 \\
        p(\text{Normal}) &= 0{,}75, \quad p(\text{Tidak normal}) = 0{,}25
        \end{aligned}
        $$
    
        \item Hitung total nilai MI dari peluang di atas:
        $$
        \begin{aligned}
        MI = & \left(0{,}25 \cdot \log_2 \frac{0{,}25}{0{,}40 \cdot 0{,}75}\right) + \left(0{,}15 \cdot \log_2 \frac{0{,}15}{0{,}40 \cdot 0{,}25}\right) + \\
             & \left(0{,}35 \cdot \log_2 \frac{0{,}35}{0{,}40 \cdot 0{,}75}\right) + \left(0{,}05 \cdot \log_2 \frac{0{,}05}{0{,}40 \cdot 0{,}25}\right) + \\
             & \left(0{,}15 \cdot \log_2 \frac{0{,}15}{0{,}20 \cdot 0{,}75}\right) + \left(0{,}05 \cdot \log_2 \frac{0{,}05}{0{,}20 \cdot 0{,}25}\right)
        \end{aligned}
        $$
        $$
        MI \approx -0{,}065 + 0{,}087 + 0{,}077 - 0{,}050 + 0 + 0 \approx 0{,}049
        $$
    \end{enumerate}
    
    Nilai MI sebesar 0,049 menunjukkan bahwa fitur lama menginap hanya memberikan sedikit informasi tambahan untuk menebak status penyewa pada kelompok data yang sangat timpang ini.

    


% Sub-bab 2.7 decision tree
\section{\it Decision Tree}

{\it Decision Tree} merupakan salah satu algoritma klasifikasi yang bekerja dengan memetakan data ke dalam kelas-kelas tertentu menggunakan urutan aturan keputusan berbentuk pola ``jika-maka'' (\textit{if-then}) (Tan, Steinbach, \& Kumar, 2019). Algoritma ini memilah kumpulan data (\textit{dataset}) menjadi kelompok-kelompok yang lebih spesifik dan seragam berdasarkan fitur (kolom) yang paling efektif dalam membedakan target kelas. Proses pemilahan ini dilakukan secara berulang (\textit{rekursif}) hingga memenuhi kriteria penghentian, misalnya saat seluruh data di dalam satu ujung cabang (\textit{node}) sudah memiliki label kelas yang sama, atau jumlah data terlalu sedikit untuk dipilah kembali (Han, Pei, \& Kamber, 2011).

Pada setiap percabangan, algoritma memilih fitur terbaik dengan mengukur tingkat "ketidakmurnian" atau ketidakpastian (\textit{impurity}) data, menggunakan metode seperti \textit{Information Gain}, \textit{Gini Index}, atau \textit{Gain Ratio}. Khusus untuk \textit{Information Gain}, rumus yang digunakan adalah sebagai berikut:
$$
IG(S, A) = Entropy(S) - \sum_{v \in Values(A)} \frac{|S_v|}{|S|} \cdot Entropy(S_v)
$$  
dengan $Entropy(S) = - \sum_{i=1}^{c} p_i \log_2 p_i$, di mana $p_i$ merupakan persentase proporsi data dalam suatu kelas. Entropi berfungsi untuk mengukur seberapa acak atau tidak pastinya sebuah himpunan data. Dengan demikian, \textit{Information Gain} pada dasarnya mengukur seberapa besar ketidakpastian tersebut berkurang setelah data dipilah menggunakan fitur $A$ (Quinlan, 1986).

Sebagai ilustrasi, simulasi perhitungan berikut diadaptasi dari \textit{dataset} klasik pada buku rujukan standar \textit{Machine Learning} (Mitchell, 1997) yang angka-angkanya dipetakan ke dalam konteks pengawasan penyewa. Misalkan terdapat 14 data penyewa ($S$) dengan fitur ``Lama Menginap'' dan ``Metode Pembayaran'', serta label kelas yang menandakan apakah penyewa tersebut tergolong normal (0) atau tidak normal (1). Mengingat pada studi kasus ini jumlah kelas tidak normal adalah minoritas, diasumsikan himpunan data terdiri atas 9 penyewa kelas normal (0) dan 5 penyewa kelas tidak normal (1). Pemilihan fitur mana yang paling cocok untuk membelah data pertama kali dilakukan dengan menghitung nilai \textit{Information Gain} (IG).

\begin{itemize}
    \item \textbf{Menghitung ketidakpastian awal (Entropi $S$)} \\
    Entropi awal dari keseluruhan himpunan 14 data penyewa dihitung sebagai berikut.
    $$
    Entropy(S) = - \left(\frac{9}{14} \log_2 \frac{9}{14}\right) - \left(\frac{5}{14} \log_2 \frac{5}{14}\right) \approx 0{,}940
    $$
    
    \item \textbf{Mengevaluasi fitur ``Lama Menginap''} \\
    Jika himpunan data dipecah menggunakan lama menginap, asumsikan data terbagi menjadi dua kelompok:
    \begin{itemize}
        \item \textbf{Kelompok 1--2 malam (7 penyewa)}, terdiri dari 3 Normal (0) dan 4 Tidak Normal (1). Karena perbandingannya cukup bervariasi, ketidakpastiannya dihitung sebagai berikut:
        $$
        Entropy(1\text{--}2\text{ malam}) = - \left(\frac{3}{7} \log_2 \frac{3}{7}\right) - \left(\frac{4}{7} \log_2 \frac{4}{7}\right) \approx 0{,}985
        $$
        \item \textbf{Kelompok $\geq$ 3 malam (7 penyewa)}, didominasi oleh kelas normal, yakni 6 Normal (0) dan 1 Tidak Normal (1). Ketidakpastiannya lebih rendah:
        $$
        Entropy(\geq3\text{ malam}) = - \left(\frac{6}{7} \log_2 \frac{6}{7}\right) - \left(\frac{1}{7} \log_2 \frac{1}{7}\right) \approx 0{,}592
        $$
    \end{itemize}
    
    \item \textbf{Menghitung \textit{Information Gain} (Lama Menginap)} \\
    Pengurangan ketidakpastian dihitung dengan mengurangkan entropi awal dengan rata-rata tertimbang (\textit{weighted average}) dari kelompok yang baru terbentuk:
    $$
    IG(S, Lama) = Entropy(S) - \left( \frac{|1\text{--}2\text{ malam}|}{|S|} \times Entropy(1\text{--}2\text{ malam}) \right) - \left( \frac{|\geq3\text{ malam}|}{|S|} \times Entropy(\geq3\text{ malam}) \right)
    $$
    $$
    IG(S, Lama) = 0{,}940 - \left( \frac{7}{14} \times 0{,}985 \right) - \left( \frac{7}{14} \times 0{,}592 \right) = 0{,}151
    $$
\end{itemize}

Setelah mendapatkan nilai evaluasi untuk Lama Menginap, perhitungan serupa dilakukan untuk fitur kandidat lainnya, yaitu ``Metode Pembayaran''. Asumsikan fitur ini memecah data menjadi kelompok ``Tunai'' dan ``Kartu''.

\begin{itemize}
    \item \textbf{Mengevaluasi fitur ``Metode Pembayaran''} \\
    Jika data dipecah berdasarkan metode pembayaran, distribusinya menjadi:
    \begin{itemize}
        \item \textbf{Kelompok Tunai (8 penyewa)}, terdiri dari 6 Normal (0) dan 2 Tidak Normal (1). Ketidakpastiannya adalah:
        $$
        Entropy(Tunai) = - \left(\frac{6}{8} \log_2 \frac{6}{8}\right) - \left(\frac{2}{8} \log_2 \frac{2}{8}\right) \approx 0{,}811
        $$
        
        \item \textbf{Kelompok Kartu (6 penyewa)}, terdiri dari 3 Normal (0) dan 3 Tidak Normal (1). Karena seimbang sempurna, ketidakpastiannya maksimal:
        $$
        Entropy(Kartu) = - \left(\frac{3}{6} \log_2 \frac{3}{6}\right) - \left(\frac{3}{6} \log_2 \frac{3}{6}\right) = 1
        $$
    \end{itemize}
    
    \item \textbf{Menghitung \textit{Information Gain} (Metode Pembayaran)} \\
    Sama seperti sebelumnya, pengurangan ketidakpastian untuk fitur ini adalah:
    $$
    IG(S, Metode) = Entropy(S) - \left( \frac{|\text{Tunai}|}{|S|} \times Entropy(Tunai) \right) - \left( \frac{|\text{Kartu}|}{|S|} \times Entropy(Kartu) \right)
    $$
    $$
    IG(S, Metode) = 0{,}940 - \left( \frac{8}{14} \times 0{,}811 \right) - \left( \frac{6}{14} \times 1 \right)
    $$
    $$
    IG(S, Metode) = 0{,}940 - 0{,}463 - 0{,}428 = 0{,}049
    $$
\end{itemize}

Berdasarkan perbandingan matematis di atas, pemisahan menggunakan fitur ``Lama Menginap'' berhasil mengurangi ketidakpastian sebesar 0,151, sedangkan ``Metode Pembayaran'' hanya 0,049 (Mitchell, 1997). Karena ``Lama Menginap'' memberikan pengurangan acak yang lebih baik ($0{,}151 > 0{,}049$), fitur tersebut secara otomatis dipilih oleh algoritma sebagai titik pembagi pertama (\textit{root node}) pada struktur pohon keputusan. Setelah titik pembagi awal ditetapkan, pohon akan terus mempercabangkan data. Penting untuk dicatat bahwa pada proses ini, algoritma tidak membuang kolom apa pun secara otomatis. Sistem akan memanfaatkan seluruh informasi relevan yang telah dikonversi (\textit{encoding}) dan mengevaluasi setiap kolom tersebut pada percabangan-percabangan selanjutnya hingga seluruh cabang bermuara pada kesimpulan akhir (daun/\textit{leaf}). Hasil akhir dari \textit{Decision Tree} adalah prediksi kelas untuk penyewa baru berdasarkan jalur kondisi yang telah dilewatinya. Semakin murni kondisi di ujung daun (misalnya berisi 100\% data tidak normal), semakin kuat keyakinan model terhadap prediksinya. Dalam penerapannya, model umumnya mengandalkan persentase kelas di ujung cabang sebagai nilai probabilitas. Sebagai contoh, jika suatu daun di ujung pohon memuat 8 data historis yang terdiri atas 6 kelas normal dan 2 kelas tidak normal, maka peluang probabilitas penyewa baru yang masuk ke daun tersebut untuk menjadi normal adalah 75\% (6/8). Angka probabilitas ini kemudian dapat dijadikan pijakan (\textit{threshold}) bagi pihak manajemen untuk menentukan tindakan; misalnya membunyikan alarm waspada jika peluang tidak normal menyentuh angka lebih dari 30\%.

Kelebihan utama dari algoritma \textit{Decision Tree} adalah kemampuannya memvisualisasikan logika pengambilan keputusan. Karena bentuknya menyerupai bagan alir dan cara kerjanya sangat identik dengan proses penalaran manusia, hasil dari model ini sangat mudah dipahami dan dijelaskan kepada pihak operasional maupun manajemen (Loh, 2011).

% Sub-bab 2.8 knn

\section{{K-Nearest Neighbors} (KNN)}

Algoritma {K-Nearest Neighbors} (KNN) merupakan metode klasifikasi non-parametrik berbasis kemiripan yang memprediksi kelas dari suatu data baru berdasarkan mayoritas kelas dari $k$ tetangga terdekatnya dalam ruang fitur. Menurut Cover dan Hart (1967) dalam jurnalnya, serta Duda, Hart, dan Stork (2001) dalam buku \textit{Pattern Classification}, KNN termasuk dalam kategori \textit{lazy learning}. Algoritma ini tidak membangun model komputasi secara eksplisit selama fase pelatihan, melainkan menyimpan seluruh data pelatihan dan menggunakan metrik jarak untuk menentukan kedekatan antar data pada saat fase pengujian. Algoritma ini sangat efektif digunakan pada data dengan distribusi yang lugas dan memiliki hubungan lokal yang kuat antar fiturnya.

\subsection{Metrik Jarak pada KNN}
Untuk mengukur seberapa ``dekat'' atau ``mirip'' suatu data baru dengan data pelatihan, KNN membutuhkan metrik jarak. Tiga jenis metrik jarak yang paling umum digunakan dalam literatur pembelajaran mesin adalah sebagai berikut.

\begin{enumerate}
    \item \textbf{Jarak Euclidean} \\
    Jarak Euclidean mengukur panjang garis lurus antara dua titik. Metode ini merupakan metrik standar dalam KNN dan sangat sesuai untuk data kontinu dengan distribusi yang proporsional dan linier.
    \[
    d(x, y) = \sqrt{\sum_{i=1}^{n} (x_i - y_i)^2}
    \]
    
    \item \textbf{Jarak Manhattan (\textit{L1 Distance})} \\
    Jarak Manhattan mengukur jumlah total selisih absolut dari tiap dimensi. Metrik ini sering digunakan ketika fitur data bersifat ordinal, diskret, atau memiliki skala yang tidak proporsional (misalnya, rentang jumlah hari yang sangat berbeda dengan kategori waktu jam).
    \[
    d(x, y) = \sum_{i=1}^{n} |x_i - y_i|
    \]
    
    \item \textbf{Jarak Minkowski} \\
    Minkowski adalah generalisasi matematis dari jarak Euclidean dan Manhattan yang dikontrol oleh parameter $p$. Jika parameter $p=1$, rumusnya menjadi jarak Manhattan, sedangkan jika $p=2$, rumusnya menjadi jarak Euclidean. Fleksibilitas parameter $p$ ini memungkinkan metrik disesuaikan dengan karakteristik unik dari sebuah dataset.
    \[
    d(x, y) = \left( \sum_{i=1}^{n} |x_i - y_i|^p \right)^{1/p}
    \]
\end{enumerate}

\subsection{Studi Kasus: Klasifikasi Status Penyewa}
Sebagai ilustrasi penerapan KNN, misalkan terdapat data penyewa dengan dua fitur, yaitu \textbf{lama menginap} (hari) dan \textbf{waktu \textit{check-in}} (jam), serta label kelas klasifikasi ``normal'' atau ``mencurigakan''. Terdapat lima data latih (A--E) sebagai berikut.

\begin{table}[H]
    \centering
    \caption{Data Latih Penyewa Properti}
    \begin{tabular}{|c|c|c|l|}
        \hline
        \textbf{Titik} & \textbf{Lama Menginap} & \textbf{Waktu \textit{Check-in}} & \textbf{Label Kelas} \\
        \hline
        A & 2 & 9 & normal \\
        B & 3 & 7 & normal \\
        C & 4 & 6 & mencurigakan \\
        D & 5 & 5 & mencurigakan \\
        E & 6 & 4 & mencurigakan \\
        \hline
    \end{tabular}
\end{table}

Misalkan terdapat sebuah data uji baru, yaitu $X(3, 6)$ yang dapat diukur kedekatannya terhadap seluruh data latih menggunakan ketiga metrik jarak di atas.

\subsubsection*{1. Perhitungan Menggunakan Jarak Euclidean}
\begin{itemize}
    \item \textbf{Jarak ke A}: $\sqrt{(3 - 2)^2 + (6 - 9)^2} = \sqrt{1 + 9} = \sqrt{10} \approx 3{,}16$
    \item \textbf{Jarak ke B}: $\sqrt{(3 - 3)^2 + (6 - 7)^2} = \sqrt{0 + 1} = \sqrt{1} = 1{,}00$
    \item \textbf{Jarak ke C}: $\sqrt{(3 - 4)^2 + (6 - 6)^2} = \sqrt{1 + 0} = \sqrt{1} = 1{,}00$
    \item \textbf{Jarak ke D}: $\sqrt{(3 - 5)^2 + (6 - 5)^2} = \sqrt{4 + 1} = \sqrt{5} \approx 2{,}24$
    \item \textbf{Jarak ke E}: $\sqrt{(3 - 6)^2 + (6 - 4)^2} = \sqrt{9 + 4} = \sqrt{13} \approx 3{,}61$
\end{itemize}

\subsubsection*{2. Perhitungan Menggunakan Jarak Manhattan}
\begin{itemize}
    \item \textbf{Jarak ke A}: $|3 - 2| + |6 - 9| = 1 + |-3| = 1 + 3 = 4$
    \item \textbf{Jarak ke B}: $|3 - 3| + |6 - 7| = 0 + |-1| = 0 + 1 = 1$
    \item \textbf{Jarak ke C}: $|3 - 4| + |6 - 6| = |-1| + 0 = 1 + 0 = 1$
    \item \textbf{Jarak ke D}: $|3 - 5| + |6 - 5| = |-2| + 1 = 2 + 1 = 3$
    \item \textbf{Jarak ke E}: $|3 - 6| + |6 - 4| = |-3| + 2 = 3 + 2 = 5$
\end{itemize}

\subsubsection*{3. Perhitungan Menggunakan Jarak Minkowski ($p=3$)}
\begin{itemize}
    \item \textbf{Jarak ke A}: $(|3 - 2|^3 + |6 - 9|^3)^{1/3} = (1 + 27)^{1/3} = \sqrt[3]{28} \approx 3{,}04$
    \item \textbf{Jarak ke B}: $(|3 - 3|^3 + |6 - 7|^3)^{1/3} = (0 + 1)^{1/3} = \sqrt[3]{1} = 1{,}00$
    \item \textbf{Jarak ke C}: $(|3 - 4|^3 + |6 - 6|^3)^{1/3} = (1 + 0)^{1/3} = \sqrt[3]{1} = 1{,}00$
    \item \textbf{Jarak ke D}: $(|3 - 5|^3 + |6 - 5|^3)^{1/3} = (8 + 1)^{1/3} = \sqrt[3]{9} \approx 2{,}08$
    \item \textbf{Jarak ke E}: $(|3 - 6|^3 + |6 - 4|^3)^{1/3} = (27 + 8)^{1/3} = \sqrt[3]{35} \approx 3{,}27$
\end{itemize}

\subsection{Penentuan Kelas Berdasarkan Metrik Jarak}
Dari hasil perhitungan matematis di atas, nilai jarak bisa bervariasi bergantung pada metrik yang digunakan. Setelah semua jarak dihitung, algoritma KNN akan mengurutkan titik-titik tersebut dari yang terdekat hingga terjauh untuk memilih $k$ tetangga. Berikut adalah evaluasi penentuan kelas untuk masing-masing metode dengan parameter $k = 3$.

\begin{itemize}
    \item \textbf{Berdasarkan Euclidean:} Urutan titik terdekat adalah \textbf{B} ($1{,}00$), \textbf{C} ($1{,}00$), \textbf{D} ($2{,}24$), A ($3{,}16$), dan E ($3{,}61$). 
    \item \textbf{Berdasarkan Manhattan:} Urutan titik terdekat adalah \textbf{B} ($1$), \textbf{C} ($1$), \textbf{D} ($3$), A ($4$), dan E ($5$).
    \item \textbf{Berdasarkan Minkowski ($p=3$):} Urutan titik terdekat adalah \textbf{B} ($1{,}00$), \textbf{C} ($1{,}00$), \textbf{D} ($2{,}08$), A ($3{,}04$), dan E ($3{,}27$).
\end{itemize}

Berdasarkan ketiga metrik tersebut, jika $k=3$, maka himpunan tetangga terdekat secara konsisten jatuh pada titik B (normal), C (mencurigakan), dan D (mencurigakan). Melalui mekanisme pemungutan suara mayoritas (\textit{majority voting}), yaitu 2 berbanding 1, maka data uji $X(3, 6)$ akan diklasifikasikan sebagai penyewa yang \textbf{mencurigakan}.

\subsection{Kesimpulan Probabilistik dan Evaluasi Algoritma}
Pada versi probabilistik, proporsi mayoritas tersebut dapat diterjemahkan menjadi probabilitas kelas. Karena 2 dari 3 tetangga terdekat berlabel ``mencurigakan'', maka skor probabilitas anomali untuk data $X$ adalah $\frac{2}{3} \approx 66{,}7\%$. Dalam implementasi sistem deteksi, dapat diterapkan ambang batas (\textit{threshold}) klasifikasi—misalnya $\geq 60\%$—untuk secara otomatis memberikan notifikasi kepada pengelola properti apabila terdapat penyewa yang terindikasi melanggar aturan. Menurut Tan, Steinbach, dan Kumar (2019) dalam buku \textit{Introduction to Data Mining}, keunggulan utama KNN terletak pada kesederhanaannya, sifatnya yang intuitif, serta kemudahan dalam interpretasi hasil prediksi. Meskipun demikian, algoritma ini memiliki kelemahan yang sangat rentan terhadap skala fitur (\textit{scale-dependent}) dan menuntut beban komputasi yang tinggi ketika ukuran data latih semakin besar. Oleh karena itu, tahapan pra-pemrosesan seperti normalisasi fitur (untuk menyamakan skala antara lama menginap dan waktu jam) serta pemilihan metrik jarak yang tepat menjadi syarat krusial agar performa KNN tetap optimal di dunia nyata.

% Sub-bab 2.9 random forest
\section{{\it Random Forest}}

{\it Random Forest} merupakan algoritma klasifikasi berbasis metode \textit{ensemble learning} yang menggabungkan hasil dari banyak pohon keputusan (\textit{decision tree}) untuk menghasilkan prediksi yang lebih stabil dan akurat. Algoritma ini dikembangkan oleh Breiman (2001) dengan tujuan meningkatkan performa klasifikasi dibandingkan pohon keputusan tunggal yang rentan terhadap {\it overfitting}. Setiap pohon dalam {\it Random Forest} dilatih menggunakan teknik {\it bootstrap aggregating} atau {\it bagging}, yaitu dengan mengambil sampel acak dari data pelatihan dengan pengembalian, serta memilih subset acak dari fitur saat membangun node pohon. Hasil akhir prediksi diperoleh dengan menggunakan mayoritas suara (\textit{majority voting}) pada klasifikasi atau rata-rata (\textit{mean}) pada regresi (Hastie, Tibshirani, \& Friedman, 2009).
    
    Secara matematis, untuk klasifikasi hasil prediksi dari {\it Random Forest} dapat diformulasikan sebagai berikut.  
    \[
    \hat{y} = \text{mode}\{h_1(x), h_2(x), ..., h_k(x)\}
    \]  
    di mana \( h_i(x) \) adalah hasil prediksi dari pohon ke-\( i \), dan \( k \) adalah jumlah total pohon. Jika mayoritas pohon memprediksi bahwa seorang penyewa termasuk kategori mencurigakan, maka model akan memberikan label mencurigakan. Selain itu, {\it Random Forest} juga dapat menghasilkan nilai probabilitas, yaitu persentase jumlah pohon yang memutuskan kelas tertentu. Sebagai contoh, jika 80 dari 100 pohon memutuskan bahwa penyewa mencurigakan, maka probabilitas penyewa mencurigakan adalah 0.80 atau 80\%.
    
    Sebagai ilustrasi, misalkan terdapat data penyewa dengan fitur lama menginap, metode pembayaran, dan waktu check-in. Data tersebut dimasukkan ke dalam model {\it Random Forest} dengan 100 pohon. Setelah proses prediksi, sebanyak 72 pohon memutuskan penyewa adalah normal dan 28 pohon memutuskan mencurigakan. Maka, probabilitas penyewa diklasifikasikan sebagai mencurigakan sebesar 28\% dan sebagai normal sebesar 72\%. Hasil ini dapat diinterpretasikan bahwa model lebih cenderung mengklasifikasikan penyewa tersebut sebagai normal, meskipun tingkat keyakinannya belum 100\%. Dalam praktik pengawasan, ambang batas (\textit{threshold}) tertentu dapat ditetapkan untuk mengklasifikasikan penyewa sebagai mencurigakan, misalnya jika probabilitas $\geq$60\%, maka akan muncul notifikasi peringatan. 
    
    Keunggulan {\it Random Forest} terletak pada kemampuannya menangani data berdimensi tinggi dan kelas yang tidak seimbang tanpa memerlukan banyak praproses. Selain itu, algoritma ini juga memberikan estimasi pentingnya fitur (\textit{feature importance}) yang bermanfaat untuk interpretasi model dan pengambilan keputusan. Dengan keunggulan tersebut, {\it Random Forest} telah digunakan secara luas dalam deteksi anomali, penipuan, dan klasifikasi risiko dalam berbagai domain, termasuk keamanan siber, kesehatan, dan properti (Liaw \& Wiener, 2002).

% Sub-bab 2.10 svm
\section{{\it Support Vector Machine} (SVM)}

{\it Support Vector Machine} (SVM) merupakan algoritma klasifikasi yang bertugas mencari batas pemisah ({\it hyperplane}) paling optimal untuk membedakan dua kelompok data. SVM bekerja dengan mencari garis (dalam dua dimensi) atau bidang (dalam dimensi lebih tinggi) yang memaksimalkan jarak antara dua kelas data yang berbeda (Cortes \& Vapnik, 1995). Dalam praktiknya, SVM efektif digunakan untuk klasifikasi data linier maupun nonlinier melalui fungsi kernel, seperti \textit{Radial Basis Function} (RBF), polinomial, dan sigmoid (Hastie, Tibshirani, \& Friedman, 2009).
    
    Secara matematis, penentuan batas pemisah pada algoritma SVM dirumuskan sebagai berikut.
    $$
    f(x) = \text{sign}(w^T x + b)
    $$
    Dalam rumus ini, $w$ merupakan bobot yang menentukan seberapa penting suatu informasi, $x$ adalah data atau fitur yang dimasukkan, dan $b$ adalah nilai penyesuaian (bias). Tujuan utama model SVM adalah mencari kombinasi nilai $w$ dan $b$ yang menghasilkan jarak pemisah paling lebar antar kelas, yang dirumuskan:
    $$
    \min \frac{1}{2} \|w\|^2 \quad \text{dengan syarat } y_i (w^T x_i + b) \geq 1
    $$
    untuk setiap data $i$. Di sini, $y_i \in \{-1, +1\}$ adalah penanda kelas data (dalam penelitian ini, nilai $-1$ berarti kelas normal (0) dan $+1$ berarti kelas tidak normal (1)). 
    
    Namun, di dunia nyata sering kali ada data yang posisinya tumpang tindih dan tidak bisa dipisahkan dengan garis batas yang sempurna. Untuk mengatasi ini, SVM menggunakan konsep batas toleransi ({\it soft margin}) dengan perhitungan kerugian sebagai berikut.
    $$
    \min \frac{1}{2} \|w\|^2 + C \sum_{i=1}^{n} \xi_i
    $$
    Simbol $\xi_i$ mewakili tingkat kelonggaran ({\it slack variable}) yang memperbolehkan sistem sesekali melakukan toleransi kesalahan batas pemisahan, sedangkan $C$ adalah parameter pengatur ketatnya toleransi tersebut. Mengingat pada pengawasan penyewa ini data kelas tidak normal (1) jumlahnya sangat sedikit dibandingkan penyewa normal (0), model rentan ``menebak aman'' dengan menganggap semua orang adalah penyewa normal. Untuk mencegah bias tersebut, SVM menerapkan sistem denda bersyarat (\textit{cost-sensitive learning}). Parameter batas toleransi $C$ dipecah sedemikian rupa sehingga jika sistem salah mengenali penyewa bermasalah (kelas minoritas), denda penalti yang diberikan ($C_1$) akan jauh lebih berat dibandingkan jika salah mengenali penyewa normal ($C_0$).
    
    Sebagai ilustrasi, simulasi penentuan batas keputusan berikut diadaptasi dari formulasi \textit{optimal separating hyperplane} yang dijabarkan dalam literatur \textit{The Elements of Statistical Learning} (Hastie, Tibshirani, \& Friedman, 2009), yang kemudian diaplikasikan pada konteks data penyewa. Misalkan terdapat dua fitur penyewa yang dinilai: lama menginap ($x_1$) dan waktu \textit{check-in} yang sudah diubah menjadi angka ($x_2$, misalnya pagi = 0, sore = 1, malam = 2). Setelah model dilatih dan bobot kepentingannya ($w$) serta nilai biasnya ($b$) disesuaikan oleh algoritma, sistem menemukan rumus keputusan batas sebagai berikut.
    $$
    f(x) = 2x_1 + 3x_2 - 5
    $$
    Misalkan ada penyewa baru yang menginap selama 3 hari ($x_1 = 3$) dan \textit{check-in} pada malam hari ($x_2 = 2$). Jika nilai tersebut dimasukkan ke dalam rumus linier di atas, perhitungannya menjadi:
    $$
    f(x) = 2(3) + 3(2) - 5 = 6 + 6 - 5 = 7
    $$
    Karena hasil akhirnya adalah angka positif (lebih besar dari 0), model akan menyimpulkan bahwa penyewa ini berada di wilayah kelas tidak normal (1). Sebaliknya, jika hasilnya berupa angka negatif ($< 0$), penyewa akan dianggap normal (0). Semakin besar angka yang dihasilkan (menjauhi nol), semakin yakin sistem terhadap tebakannya. Jika angkanya sangat mendekati nol, berarti posisi penyewa tersebut berada tepat di pinggir garis batas keraguan.
    
    Pada kenyataannya, banyak data di dunia nyata memiliki pola perilaku yang rumit sehingga tidak bisa dipisahkan hanya dengan sebuah garis pemisah yang lurus (linier). Untuk mengatasi hambatan tersebut, SVM mengandalkan sebuah strategi matematis berjuluk ``Trik Kernel'' (\textit{Kernel Trick}). Fungsi kernel memanipulasi rentang batas fitur yang saling bertumpuk di ruang dimensi asalnya, memproyeksikannya secara otomatis ke sebuah ruang khayal berdimensi jauh lebih tinggi, sehingga dapat diiris dengan sempurna membentang lurus oleh batas keputusan (Sch\"{o}lkopf \& Smola, 2002).

    Salah satu fungsi kernel yang paling dapat diandalkan dan luas penerapannya adalah \textit{Radial Basis Function} (RBF). Secara matematis, tingkat kemiripan atau keterkaitan (\textit{similarity}) antara dua buah titik data awal ($x_i$ dan $x_j$) menggunakan fungsi RBF dirumuskan dalam persamaan berikut (Hastie, Tibshirani, \& Friedman, 2009).
    $$
    K(x_i, x_j) = \exp(-\gamma \|x_i - x_j\|^2)
    $$
    Pada rumus tersebut, simbol $\gamma$ (gamma) berfungsi sebagai pengatur seberapa ketat atau melengkung batas pengelompokan yang akan dibuat, sedangkan simbol $\|x_i - x_j\|^2$ menunjukkan jarak lurus (\textit{Euclidean distance}) antara dua titik data.  Secara sederhana, apabila dua data memiliki karakteristik yang sangat mirip atau posisinya saling berdekatan, maka hasil perhitungan dari rumus ini akan bernilai mendekati angka 1 ($\exp(0) = 1$). Sebaliknya, jika kedua data tersebut sangat berbeda atau posisinya saling berjauhan, maka nilainya akan mendekati angka 0. Kemampuan RBF dalam membungkus data berbentuk sebaran kurva geometris melingkar, menjadikannya senjata mumpuni bagi SVM saat memilah distribusi sebaran data rentan (\textit{non-linear}) tanpa memaksa para analis repot memikirkan rekayasa rumusnya secara manual (Burges, 1998).
    
    Algoritma SVM dikenal sangat tangguh untuk memproses data yang memiliki banyak kolom fitur. Selain itu, SVM juga cukup kebal terhadap kecenderungan menghafal data secara berlebihan ({\it overfitting}), terutama pada dataset yang minim gangguan ({\it noise}). Karena sifatnya yang mendefinisikan batas keputusan eksplisit, SVM banyak digunakan dalam bidang deteksi intrusi, identifikasi penipuan, dan klasifikasi risiko, termasuk dalam pengawasan properti berbasis data (Sch\"{o}lkopf \& Smola, 2002).

% Sub-bab 2.11 xgboost
\section{\it XGBoost}

{\it Extreme Gradient Boosting} (XGBoost) merupakan algoritma klasifikasi berbasis pohon keputusan yang dikembangkan dari metode {\it Gradient Boosting Machine} (GBM) dan dioptimalkan untuk efisiensi, skalabilitas, serta akurasi (Chen \& Guestrin, 2016). XGBoost bekerja dengan cara membangun model secara bertahap melalui penambahan pohon-pohon keputusan secara berurutan untuk memperbaiki kesalahan prediksi dari model sebelumnya. Pada setiap iterasi, model baru dilatih untuk meminimalkan fungsi kerugian yang dikombinasikan dengan regularisasi guna mencegah {\it overfitting} (Friedman, 2001).

    Fungsi objektif pada XGBoost dirumuskan sebagai berikut:
    \[
    \mathcal{L}(\phi) = \sum_{i=1}^{n} l(\hat{y}_i, y_i) + \sum_{k=1}^{K} \Omega(f_k)
    \]
    dengan \(\Omega(f) = \gamma T + \frac{1}{2} \lambda \sum_{j=1}^{T} w_j^2\),  
    di mana \(l\) adalah fungsi kerugian (misalnya log-loss), \(\Omega\) adalah fungsi regularisasi, dan \(T\) adalah jumlah daun pada pohon, dan \(w_j\) bobot pada daun ke-$j$. Tujuan dari XGBoost adalah meminimalkan fungsi \(\mathcal{L}\), sehingga model menjadi akurat dan tidak terlalu kompleks.
    
    Misalkan terdapat tiga data penyewa untuk pelatihan:
    \begin{itemize}
        \item Data 1: (lama menginap = 2 hari, check-in = 22.00) $\rightarrow$ label: 0 (normal)
        \item Data 2: (lama menginap = 1 hari, check-in = 00.00) $\rightarrow$ label: 1 (mencurigakan)
        \item Data 3: (lama menginap = 3 hari, check-in = 20.00) $\rightarrow$ label: 0 (normal)
    \end{itemize}
    
    Asumsikan prediksi awal (\(\hat{y}_i\)) untuk semua data adalah 0.5, maka residu dihitung menggunakan derivatif pertama dari log-loss.
    \[
    g_i = \hat{y}_i - y_i
    \] maka: 
    \[
    \text{Data 1: } g_1 = 0{.}5 - 0 = 0{.}5,\quad 
    \text{Data 2: } g_2 = 0{.}5 - 1 = -0{.}5,\quad 
    \text{Data 3: } g_3 = 0{.}5 - 0 = 0{.}5
    \]
    
    Kemudian pohon pertama dibentuk berdasarkan {\it gradient} tersebut. Misalkan pohon tersebut membagi data berdasarkan \textit{check-in}:
    \begin{itemize}
      \item jika \textit{check-in} $>$ 21.00 $\rightarrow$ ke daun A
      \item jika \textit{check-in} $\leq$ 21.00 $\rightarrow$ ke daun B
    \end{itemize}
    
    maka:
    \begin{itemize}
      \item data 1 dan 2 masuk daun A
      \item data 3 masuk daun B
    \end{itemize}
    
    Langkah selanjutnya adalah menghitung nilai bobot prediksi di tiap daun menggunakan rumus berikut.
    \[
    w_j = -\frac{\sum g_i}{\sum h_i + \lambda}
    \]
    Dengan \(g_i\) adalah {\it gradient} dan \(h_i = \hat{y}_i (1 - \hat{y}_i)\) adalah {\it hessian} (turunan kedua dari log-loss). Misal \(\lambda = 1\), dan semua \(h_i = 0{,}5 \cdot 0{,}5 = 0{,}25\).
    
    Untuk daun A (Data 1 dan 2):
    \[
    \sum g_i = 0{.}5 + (-0{.}5) = 0,\quad
    \sum h_i = 0{.}25 + 0{.}25 = 0{.}5
    \Rightarrow w_A = -\frac{0}{0{.}5 + 1} = 0
    \]
    
    Untuk daun B (Data 3):
    \[
    \sum g_i = 0{.}5,\quad
    \sum h_i = 0{.}25
    \Rightarrow w_B = -\frac{0{.}5}{0{.}25 + 1} = -\frac{0{.}5}{1{.}25} = -0{.}4
    \]
    % \newpage
    Pohon ini kemudian akan memberikan nilai output:
    \begin{itemize}
      \item Untuk daun A: $0$
      \item Untuk daun B: $-0{.}4$
    \end{itemize}
    
    
    Output tersebut dijumlahkan dengan prediksi sebelumnya (0{.}5) untuk mendapatkan prediksi baru menggunakan fungsi logit:
    \[
    \hat{y}_{baru} = \sigma(\hat{y}_{lama} + \eta \cdot w_j)
    \]
    Dengan \(\eta =\) {\it learning rate} (misal 0{.}3) dan \(\sigma\) adalah sigmoid:
    \[
    \sigma(z) = \frac{1}{1 + e^{-z}}
    \]
    
    Untuk Data 3:
    \[
    z = 0{,}5 + (0{,}3 \cdot -0{,}4) = 0{,}5 - 0{,}12 = 0{,}38
    \Rightarrow \hat{y} = \frac{1}{1 + e^{-0{,}38}} \approx 0{,}594
    \]
    
    Dikarenakan prediksi masih di bawah batas klasifikasi 0{.}60, maka penyewa diklasifikasikan sebagai normal. Jika nilai akhir $\geq$ 0{.}60, maka akan diklasifikasikan sebagai mencurigakan. Model akan melanjutkan iterasi dengan membentuk pohon-pohon tambahan untuk terus memperbaiki kesalahan prediksi sebelumnya.
    
    Keunggulan XGBoost terletak pada kemampuannya menangani data tidak seimbang menggunakan parameter seperti {\tt scale\_pos\_weight}, serta mendukung evaluasi pentingnya fitur melalui teknik {\it feature importance}. Hal ini menjadikannya sangat cocok untuk aplikasi yang menuntut akurasi tinggi serta kemampuan interpretasi terhadap pola perilaku, seperti pada sistem pengawasan penyewa mencurigakan (Kou et al., 2021).

% Sub-bab 2.12 cost sensitive learning
\section{\textit{Cost-Sensitive Learning}}

\textit{Cost-Sensitive Learning} adalah metode penanganan ketidakseimbangan kelas (\textit{class imbalance}) yang bekerja dengan cara memberikan hukuman (penalti) atau bobot kesalahan yang lebih besar saat model gagal mendeteksi kelas minoritas. Berbeda dengan teknik memanipulasi jumlah data (seperti menggandakan atau menghapus baris data asli), metode ini langsung mengubah cara algoritma belajar sehingga keaslian pola data tetap terjaga (He \& Garcia, 2009).

Pada sistem klasifikasi standar, semua kesalahan dianggap memiliki tingkat kerugian yang sama—baik itu salah memprediksi penyewa normal (0) sebagai tidak normal (1), maupun sebaliknya. Namun, di dunia nyata, gagal mendeteksi penyewa yang bermasalah (\textit{false negative}) membawa risiko bahaya dan kerugian finansial yang jauh lebih besar. Sebaliknya, salah mencurigai penyewa yang sebenarnya normal (\textit{false positive}) umumnya hanya sebatas memunculkan peringatan palsu yang mudah diverifikasi (Elkan, 2001).

Secara matematis, pendekatan ini menggunakan matriks biaya (\textit{cost matrix}) $C$, di mana $C(i, j)$ adalah skor kerugian jika model memprediksi kelas $i$ padahal data aslinya adalah kelas $j$. Jika prediksi model benar, maka tidak ada kerugian yang dihitung ($C(i, i) = 0$). Perhitungan kerugian yang diharapkan (\textit{expected cost}) dari sebuah data $x$ jika diprediksi sebagai kelas $i$ dirumuskan sebagai berikut:
$$
L(x, i) = \sum_{j} P(j \mid x) C(i, j)
$$
dengan $P(j \mid x)$ adalah probabilitas (peluang) bahwa data $x$ sesungguhnya adalah kelas $j$ berdasarkan perhitungan model. Model kemudian akan memilih prediksi kelas yang menghasilkan total kerugian $L(x, i)$ paling kecil.

Sebagai contoh perhitungan, misalkan ditetapkan matriks biaya sebagai berikut:
\begin{itemize}
    \item Biaya salah memprediksi Normal padahal aslinya Tidak Normal: $C(0, 1) = 5$ (kerugian besar).
    \item Biaya salah memprediksi Tidak Normal padahal aslinya Normal: $C(1, 0) = 1$ (kerugian kecil).
\end{itemize}

Suatu saat, model mendeteksi profil penyewa baru dan menghasilkan probabilitas dasar 70\% kemungkinan ia adalah penyewa Normal ($P(0 \mid x) = 0{,}7$) dan 30\% kemungkinan ia Tidak Normal ($P(1 \mid x) = 0{,}3$).
Model akan menghitung kerugian dari masing-masing pilihan keputusan:
\begin{itemize}
    \item Jika model memutuskan memprediksi \textbf{Normal (0)}:
    $$
    L(x, 0) = (P(0 \mid x) \times C(0, 0)) + (P(1 \mid x) \times C(0, 1))
    $$
    $$
    L(x, 0) = (0{,}7 \times 0) + (0{,}3 \times 5) = 1{,}5
    $$
    
    \item Jika model memutuskan memprediksi \textbf{Tidak Normal (1)}:
    $$
    L(x, 1) = (P(0 \mid x) \times C(1, 0)) + (P(1 \mid x) \times C(1, 1))
    $$
    $$
    L(x, 1) = (0{,}7 \times 1) + (0{,}3 \times 0) = 0{,}7
    $$
\end{itemize}

Karena kerugian memprediksi Tidak Normal lebih kecil ($0{,}7 < 1{,}5$), maka algoritma akan mengeluarkan \textit{output} \textbf{Tidak Normal (1)}. Contoh ini membuktikan bahwa meskipun secara probabilitas dasar penyewa tersebut lebih condong ke arah Normal (70\%), model tetap memilih klasifikasi Tidak Normal demi menghindari risiko fatal dari kesalahan deteksi. 

Dalam penerapannya pada program \textit{machine learning} saat ini, matriks biaya tersebut diterapkan secara praktis ke dalam bentuk rasio pembobotan kelas (\texttt{class\_weight}). Penentuan besaran angka bobot ini tidak dilakukan secara acak, melainkan dihitung secara matematis agar mencerminkan perbandingan terbalik dari jumlah distribusi masing-masing kelas target di dalam \textit{dataset} (Sun et al., 2009). Secara sistematis, formula untuk mendapatkan rasio bobot yang seimbang tersebut dirumuskan sebagai berikut.
$$
\text{Bobot Kelas} = \frac{\text{Total Label Data}}{2 \times \text{Jumlah Data pada Kelas Tersebut}}
$$

Sebagai contoh, misalkan sistem dilatih menggunakan 1.000 data riwayat penyewa, yang terdiri atas 800 data penyewa normal (mayoritas, bernilai 0) dan 200 data penyewa bermasalah (minoritas, bernilai 1). Menggunakan rumus di atas, algoritma akan menghasilkan nilai pembobotan sebesar 0,625 untuk kelas normal dan 2,5 untuk kelas bermasalah (\texttt{class\_weight=\{0:0{,}625, 1:2{,}5\}}). Secara proporsi matematis, perbandingan antara 0,625 dan 2,5 adalah selaras dengan rasio satu berbanding empat ($1:4$). Hal ini memberikan instruksi kepada model bahwa jika sistem sampai gagal mendeteksi satu saja penyewa bermasalah (\textit{false negative}), hukuman (penalti) yang diberikan pada proses perhitungannya adalah empat kali lipat lebih besar dibandingkan jika sistem keliru mencurigai penyewa normal (\textit{false positive}).

Pada beberapa algoritma modern seperti \textit{Extreme Gradient Boosting} (XGBoost), mekanisme penalti berbobot ini disederhanakan kembali ke dalam satu pengaturan parameter fungsional, yaitu skala rasio positif (\texttt{scale\_pos\_weight}). Nilai parameter ini didapatkan cukup dengan membagi total data kelas mayoritas dengan total data kelas minoritas. Keunggulan utama dari penerapan \textit{Cost-Sensitive Learning} melalui penyesuaian bobot matematis ini adalah kemurnian struktur dan pola (\textit{pattern}) asli dari \textit{dataset} tetap terjaga tanpa terdistorsi. Hal ini berbeda dengan metode penyeimbangan lain (seperti menggandakan atau memotong jumlah baris data) yang dapat merusak keaslian informasi dan sangat membebani memori komputer saat proses pelatihan berlangsung.

% Sub-bab 2.13 hyperparameter tuning
\section{\textit{Hyperparameter Tuning}}\label{sec:hyperparameter_tuning}

Dalam membangun model \textit{machine learning}, terdapat dua jenis parameter, yakni parameter internal yang dipelajari secara otomatis dari data oleh algoritma selama proses pelatihan dan \textit{hyperparameter} yang konfigurasinya harus ditentukan secara manual oleh pengembang sebelum proses pelatihan dimulai (Géron, 2019). \textit{Hyperparameter tuning} adalah tahapan kritis di mana pengembang bereksperimen mencari kombinasi \textit{hyperparameter} yang paling optimal guna menghasilkan model dengan performa prediksi terbaik dan minim tingkat kesalahan generalisasi pada data baru.

Sebuah konsep yang menjadi tujuan utama dalam proses \textit{tuning} adalah mencapai apa yang sering diistilahkan sebagai \textit{Sweet Spot}. Menurut Kuhn dan Johnson (2013) dalam \textit{Applied Predictive Modeling}, \textit{Sweet Spot} merupakan titik keseimbangan kompleksitas ideal di mana sebuah algoritma tidak terjebak pada kondisi \textit{underfitting} (terlalu sederhana hingga gagal menangkap pola) maupun \textit{overfitting} (terlalu rumit hingga menghafal gangguan atau \textit{noise} pada data latih secara presisi namun gagal saat digunakan memprediksi data nyata). Menemukan \textit{Sweet Spot} memastikan bahwa model benar-benar mempelajari pola fundamental yang relevan.

Selain itu, pemilihan arsitektur model dalam proses optimasi hyperparameter sangat dipandu oleh prinsip \textit{Parsimony} (\textit{principle of parsimony} atau \textit{Occam's Razor}). Prinsip \textit{parsimony} menyatakan bahwa jika terdapat beberapa model dengan performa prediksi yang serupa, maka model dengan struktur atau jumlah parameter yang paling sederhana dan paling transparan harus selalu diutamakan (Kuhn \& Johnson, 2013). Model yang menjunjung tinggi \textit{parsimony} jauh lebih efisien dalam hal komputasi, berjalan lebih cepat saat diimplementasikan ke dalam program nyata (\textit{deployment}), serta jauh lebih mudah dijelaskan alur cara kerjanya (\textit{interpretability}) kepada pemangku kepentingan tanpa mengorbankan akurasi yang berarti.

Salah satu metode modern yang terbukti sangat andal dan umum digunakan untuk mencari \textit{Sweet Spot} tersebut secara komputasional tanpa menguji keseluruhan probabilitas adalah \textit{Randomized Search Cross-Validation} (\texttt{RandomizedSearchCV}). Berbeda dengan metode penelusuran tradisional yang bersifat komprehensif (\textit{Grid Search}) sehingga memakan waktu ekstrem akibat permutasi kombinasi jutaan nilai, metode penelusuran acak menyelesaikan masalah tersebut dengan mengevaluasi sejumlah kombinasi acak dengan rentang iterasi dan probabilitas tertentu (Bergstra \& Bengio, 2012). Secara matematis dan probabilistik, mengambil sampel acak dalam jumlah yang memadai dari kumpulan \textit{hyperparameter} dapat menemukan performa sistem yang nyaris identik dengan hasil penelusuran \textit{grid} lengkap yang komprehensif. Pendekatan ini memungkinkan optimasi model dalam waktu yang dipercepat (Raschka \& Mirjalili, 2019).


% Sub-bab 2.14 evaluasi_model
\section{Evaluasi Model}

Evaluasi performa model merupakan tahapan penting dalam pembelajaran mesin (\textit{machine learning}) untuk menilai seberapa baik model mampu melakukan prediksi terhadap data yang belum pernah dilihat sebelumnya. Menurut Tan, Steinbach, dan Kumar (2019) dalam buku \textit{Introduction to Data Mining}, beberapa metrik umum yang digunakan dalam evaluasi model klasifikasi mencakup akurasi, presisi, recall, dan F1-score. Masing-masing metrik memberikan perspektif berbeda terhadap performa model, sehingga pemilihan metrik yang tepat bergantung pada konteks dan karakteristik data.
    
    \begin{itemize}
        \item {Akurasi} (\textit{accuracy}) adalah rasio jumlah prediksi yang benar terhadap total prediksi.
        \[
        \text{Akurasi} = \frac{TP + TN}{TP + TN + FP + FN}
        \]
        dengan:
        \begin{itemize}
            \item $TP$ (True Positive): jumlah kasus positif yang berhasil diprediksi dengan benar,
            \item $TN$ (True Negative): jumlah kasus negatif yang berhasil diprediksi dengan benar,
            \item $FP$ (False Positive): jumlah kasus negatif yang salah diklasifikasikan sebagai positif,
            \item $FN$ (False Negative): jumlah kasus positif yang salah diklasifikasikan sebagai negatif.
        \end{itemize}
        Meskipun akurasi memberikan gambaran umum performa model, metrik ini dapat menyesatkan pada dataset yang tidak seimbang. Misalnya, jika hanya 5\% dari total data merupakan penyewa mencurigakan, model yang selalu memprediksi ``normal'' tetap bisa memiliki akurasi tinggi.
    
        \item {Presisi} (\textit{precision}) mengukur ketepatan prediksi positif, yaitu seberapa besar proporsi prediksi positif yang benar-benar positif.
        \[
        \text{Presisi} = \frac{TP}{TP + FP}
        \]
        Presisi tinggi menunjukkan bahwa model jarang menghasilkan alarm palsu. Dalam konteks deteksi penyewa mencurigakan, presisi penting karena kesalahan dalam menandai penyewa normal sebagai mencurigakan dapat berdampak sosial dan hukum.
    
        \item {Recall} (atau \textit{sensitivity}) mengukur seberapa baik model dalam menemukan seluruh kasus positif.
        \[
        \text{Recall} = \frac{TP}{TP + FN}
        \]
        Recall tinggi menunjukkan bahwa model berhasil menangkap hampir semua penyewa yang benar-benar mencurigakan. Namun, jika hanya fokus pada recall, model bisa menghasilkan banyak alarm palsu sehingga presisinya rendah.
    
        \item {F1-score} adalah rata-rata harmonis dari presisi dan recall.
        \[
        \text{F1-score} = 2 \cdot \frac{\text{Presisi} \cdot \text{Recall}}{\text{Presisi} + \text{Recall}}
        \]
        Nilai F1-score berada antara 0 hingga 1. Nilai mendekati 1 menunjukkan bahwa model memiliki keseimbangan yang baik antara presisi dan recall, sehingga cocok digunakan ketika keduanya sama-sama penting.
    \end{itemize}
    
    

    Untuk menggambarkan kelemahan metrik akurasi pada dataset yang tidak seimbang, literatur standar pembelajaran mesin sering menggunakan skenario deteksi anomali, seperti penipuan kartu kredit atau deteksi penyakit langka (Provost \& Fawcett, 2013). Sebagai ilustrasi yang diadaptasi dari studi kasus evaluasi model klasifikasi, misalkan terdapat 10.000 data transaksi (dengan analogi data penyewa), di mana hanya 100 di antaranya adalah kasus kecurangan nyata (mencurigakan), sedangkan 9.900 sisanya adalah transaksi normal. Sebuah model deteksi mencurigai 150 transaksi, dan dari jumlah tersebut sebesar 90 adalah benar-benar kecurangan, maka:
    
    \begin{itemize}
        \item $TP = 90$ (terdeteksi mencurigakan dan benar kecurangan)
        \item $FP = 60$ (terdeteksi mencurigakan, tetapi sebenarnya normal)
        \item $FN = 10$ (tidak terdeteksi, padahal kecurangan)
        \item $TN = 9840$ (tidak terdeteksi dan memang normal)
    \end{itemize}
            
    Berikut perhitungannya.
            
    \[
    \text{Akurasi} = \frac{TP + TN}{10000} = \frac{90 + 9840}{10000} = 0{,}993 \quad (99{,}3\%)
    \]
    \[
    \text{Presisi} = \frac{90}{90 + 60} = \frac{90}{150} = 0{,}60 \quad (60\%)
    \]
    \[
    \text{Recall} = \frac{90}{90 + 10} = \frac{90}{100} = 0{,}90 \quad (90\%)
    \]
    \[
    \text{F1-score} = 2 \cdot \frac{0{,}60 \cdot 0{,}90}{0{,}60 + 0{,}90} = \frac{1{,}08}{1{,}50} = 0{,}72 \quad (72\%)
    \]
            
    Hasil evaluasi ini menunjukkan fenomena yang dikenal sebagai Paradoks Akurasi (\textit{Accuracy Paradox}). Meskipun model seolah-olah berkinerja sangat luar biasa dengan akurasi 99,3\%, nilai presisinya hanya 60\%. Hal ini berarti 40\% dari total peringatan yang dikeluarkan oleh sistem adalah alarm palsu (\textit{false positive}). Dalam konteks penyewaan properti, jika pengelola hanya melihat nilai akurasi, mereka mungkin menganggap sistem sudah sempurna. Namun, presisi 60\% menunjukkan risiko tinggi di mana banyak penyewa normal yang baik-baik saja akan dituduh atau ditandai sebagai penyewa mencurigakan. Kesalahan identifikasi ini dapat memicu konflik dengan penyewa, pelanggaran privasi, atau hilangnya pendapatan penyewaan. 
    
    Kelemahan metrik ini sejalan dengan temuan Fawcett (2006) dalam \textit{An Introduction to ROC Analysis}, yang menyatakan bahwa akurasi cenderung tidak dapat diandalkan ketika jumlah kasus positif sangat sedikit. Lebih lanjut, menurut Saito dan Rehmsmeier (2015) dalam jurnal \textit{PLOS ONE}, metrik presisi dan \textit{recall} (termasuk \textit{F1-score}) terbukti lebih informatif dan memberikan gambaran performa yang jauh lebih jujur dibandingkan akurasi pada dataset yang tidak seimbang. Maka dari itu, evaluasi performa model dalam penelitian ini akan lebih menitikberatkan pada metrik presisi, guna meminimalkan tingkat \textit{false positive} dan menjaga keandalan sistem dalam mendeteksi penyewa mencurigakan.





    
    

% Sub-bab 2.15 gcp
\section{Google Cloud Platform (GCP)}

Google Cloud Platform (GCP) adalah layanan komputasi awan yang dikembangkan oleh Google untuk menyediakan berbagai layanan infrastruktur, analitik data, dan kecerdasan buatan (AI). Dengan menggunakan GCP, pengguna dapat menyimpan, mengelola, dan memproses data dalam skala besar tanpa harus mengelola server fisik secara langsung. Salah satu fitur unggulan GCP adalah kemampuannya dalam mendukung pengembangan model {\it machine learning} dengan berbagai alat, seperti Google Cloud Storage untuk penyimpanan data, BigQuery untuk analitik data, dan Vertex AI untuk {\it model deployment}. Selain itu, GCP mendukung berbagai {\it framework} populer seperti {\it TensorFlow, Scikit-learn}, dan {\it PyTorch} yang memungkinkan pengguna untuk menerapkan model {\it machine learning}.

GCP memiliki beberapa keunggulan utama yang membuatnya menjadi pilihan utama dalam penerapan {\it machine learning}. Salah satu kelebihannya adalah skalabilitas dan fleksibilitas yang memungkinkan pengguna untuk menyesuaikan kapasitas penyimpanan dan daya komputasi sesuai kebutuhan. Hal ini sangat berguna ketika menangani data dalam jumlah besar atau melakukan pelatihan model dengan sumber daya komputasi yang tinggi. Selain itu, GCP menawarkan layanan {\it machine learning} terintegrasi yang mana mempermudah pengembangan, pelatihan, dan penyebaran model tanpa harus mengatur infrastruktur secara manual. Dari segi keamanan, GCP memiliki sistem keamanan tingkat tinggi, termasuk enkripsi data saat transit dan saat diam, serta kontrol akses berbasis peran untuk melindungi data pengguna. Keunggulan lainnya adalah kemampuan analitik data yang kuat di mana layanan seperti BigQuery memungkinkan analisis data dalam skala besar dengan performa tinggi. Namun, GCP juga memiliki beberapa kekurangan yang perlu diperhatikan. Biaya penggunaan GCP bisa menjadi mahal jika tidak dikonfigurasi dengan baik, terutama untuk proyek {\it machine learning} skala besar yang membutuhkan banyak sumber daya komputasi. 

Untuk mengimplementasikan model deteksi anomali seperti {\it Isolation Forest} dan {\it One-Class SVM} di GCP, dibutuhkan serangkaian tahapan mulai dari persiapan data hingga penyebaran model. Tahapan ini bersifat umum dan fleksibel untuk model-model berbasis Python dengan langkah-langkah sebagai berikut\footnote{\url{https://medium.com/publicis-sapient-france/how-to-deploy-your-own-ml-model-to-gcp-in-5-simple-steps-bf2b5898c1ab}, diakses pada 17 Maret 2025}.
    
    \begin{enumerate}
      \item {Persiapan data}
      \begin{enumerate}
        \item Mengumpulkan data menggunakan {\it data capture} yang terhubung dengan GCP. Data ini mencakup identitas penyewa, metode pembayaran, waktu {\it check-in/out}, dan lama menginap. Data disimpan di Google Cloud Storage (GCS) untuk diproses lebih lanjut.
        \item Membersihkan data menggunakan \texttt{pandas}:
        \begin{lstlisting}[language=Python]
            import pandas as pd
            # Membaca dataset dari GCS
            data = pd.read_csv("data_capture.csv")
            # Menghapus nilai kosong (NaN)
            data.dropna(inplace=True)
        \end{lstlisting}
        
        \item Melakukan normalisasi data menggunakan \texttt{StandardScaler}:
        \begin{lstlisting}[language=Python]
            from sklearn.preprocessing import StandardScaler
            scaler = StandardScaler()
            data_scaled = scaler.fit_transform(data)
        \end{lstlisting}
        
        \item Menyimpan kembali data bersih ke GCS:
        \begin{lstlisting}[language=Python]
            from google.cloud import storage
            client = storage.Client()
            bucket = client.bucket("nama-bucket")
            blob = bucket.blob("cleaned_dataset.csv")
            blob.upload_from_filename("cleaned_data.csv")
        \end{lstlisting}
      \end{enumerate}
    
      \item {Pengembangan model {\it machine learning}}
      \begin{enumerate}
        \item Mengimpor pustaka yang dibutuhkan:
        \begin{lstlisting}[language=Python]
            import numpy as np
            import pandas as pd
            from sklearn.ensemble import IsolationForest
            from sklearn.svm import OneClassSVM
            from google.cloud import storage
        \end{lstlisting}
        
        \item Membaca data dari GCS:
        \begin{lstlisting}[language=Python]
            client = storage.Client()
            bucket = client.get_bucket('nama-bucket')
            blob = bucket.blob('cleaned_dataset.csv')
            data = pd.read_csv(blob.open('r'))
        \end{lstlisting}
        
        \item Membagi data menjadi {\it training set} dan {\it testing set}:
        \begin{lstlisting}[language=Python]
            from sklearn.model_selection import train_test_split
            X_train, X_test = train_test_split(data, test_size=0.2, random_state=42)
        \end{lstlisting}
        
        \item Melatih model {\it Isolation Forest} dan {\it One-Class SVM}:
        \begin{lstlisting}[language=Python]
            # Isolation Forest
            model_if = IsolationForest(n_estimators=100, contamination=0.1, random_state=42)
            model_if.fit(X_train)
            
            # One-Class SVM
            model_svm = OneClassSVM(kernel='rbf', gamma='auto')
            model_svm.fit(X_train)
        \end{lstlisting}
      \end{enumerate}
    
      \item {Evaluasi model}
      \begin{enumerate}
        \item Melakukan prediksi pada {\it testing set}:
        \begin{lstlisting}[language=Python]
            y_pred_if = model_if.predict(X_test)
            y_pred_svm = model_svm.predict(X_test)
        \end{lstlisting}
        
        \item Mengukur performa model dengan \texttt{classification\_report}:
        \begin{lstlisting}[language=Python]
            from sklearn.metrics import classification_report
            print(classification_report(y_test, y_pred_if))
            print(classification_report(y_test, y_pred_svm))
        \end{lstlisting}
      \end{enumerate}
    
      \item {Deploy model ke Google AI Platform}
      \begin{enumerate}
        \item Menyimpan model ke dalam format \texttt{.pkl}:
        \begin{lstlisting}[language=Python]
            import joblib
            joblib.dump(model_if, 'isolation_forest.pkl')
            joblib.dump(model_svm, 'one_class_svm.pkl')
        \end{lstlisting}
        
        \item Mengunggah model ke GCS:
        \begin{lstlisting}[language=Python]
            bucket.blob('models/isolation_forest.pkl').upload_from_filename('isolation_forest.pkl')
            bucket.blob('models/one_class_svm.pkl').upload_from_filename('one_class_svm.pkl')
        \end{lstlisting}
        
        \item Melakukan deploy ke Vertex AI agar dapat digunakan untuk inferensi {\it real-time}. Sistem ini akan mendukung deteksi anomali secara otomatis dan pengambilan keputusan berbasis data.
      \end{enumerate}
\end{enumerate}

% Sub-bab 2.16 python_streamlit
% \newpage

\section{Python dan Streamlit}\label{sec:python_streamlit}

{Python} merupakan bahasa pemrograman tingkat tinggi yang sangat populer dalam pengembangan aplikasi {\it data science}, termasuk untuk implementasi {\it machine learning} dan visualisasi data interaktif. Dengan ekosistem pustaka yang kaya seperti {\it pandas, scikit-learn, matplotlib}, dan {\it seaborn}, {Python} memungkinkan pengguna untuk melakukan {\it preprocessing data}, membangun model deteksi anomali, serta melakukan evaluasi performa model secara efisien.

    Untuk kebutuhan visualisasi dan interaksi antarmuka pengguna, Streamlit menjadi solusi ringan dan efisien. Streamlit adalah {\it framework} berbasis {Python} yang memungkinkan pembuatan antarmuka web interaktif untuk aplikasi {\it machine learning} hanya dengan beberapa baris kode. Dengan Streamlit, pengguna dapat membangun sistem {\it front-end} secara cepat untuk menginput data, menampilkan hasil analisis, serta memberikan interaksi langsung seperti unggah data, filter, dan visualisasi hasil model jika diperlukan.
    
    Keunggulan utama kombinasi {Python} dan Streamlit dalam konteks penelitian ini adalah kemudahan integrasi antara logika {\it backend} dan {\it frontend} untuk interaksi pengguna. Streamlit juga mendukung proses {\it deployment} yang mudah melalui layanan seperti Streamlit Community Cloud, tanpa perlu pengaturan server manual. Ini sangat sesuai untuk sistem skala kecil hingga menengah, seperti sistem pemantauan perilaku penyewa di rusunami. Keunggulan {Python} dalam pemrosesan data dan algoritma {\it machine learning}, ditambah dengan kemampuan visualisasi interaktif dari Streamlit, membuatnya ideal digunakan dalam deteksi perilaku di hunian vertikal seperti rusunami. Berikut langkah-langkah implementasi {Python} untuk implementasi model {\it machine learning}.
    
    \begin{enumerate}
        \item {Persiapan data}
        \begin{enumerate}
            \item Mengumpulkan data penyewa yang meliputi identitas, metode pembayaran, waktu {\it check-in/out}, dan durasi menginap.
            \item Membersihkan data dengan pandas:
                \begin{lstlisting}[language=Python]
    import pandas as pd
    data = pd.read_csv("data_capture.csv")
    data.dropna(inplace=True)
                \end{lstlisting}
            \item Melakukan normalisasi data
                \begin{lstlisting}[language=Python]
    from sklearn.preprocessing import StandardScaler
    scaler = StandardScaler()
    data_scaled = scaler.fit_transform(data)
                \end{lstlisting}
        \end{enumerate}
    
        \item {Pengembangan model {\it machine learning}}
            \begin{enumerate}
                \item Import pustaka dan bagi data menjadi {\it training set} dan {\it tes set}:
                    \begin{lstlisting}[language=Python]
    from sklearn.ensemble import IsolationForest
    from sklearn.svm import OneClassSVM
    from sklearn.model_selection import train_test_split
    
    X_train, X_test = train_test_split(data_scaled, test_size=0.2, random_state=42)
                    \end{lstlisting}
        
                \item Melatih model
                    \begin{lstlisting}[language=Python]
    model_if = IsolationForest(n_estimators=100, contamination=0.1)
    model_if.fit(X_train)
    
    model_svm = OneClassSVM(kernel='rbf', gamma='auto')
    model_svm.fit(X_train)
                    \end{lstlisting}
            \end{enumerate}
    
        \item {Evaluasi model}\\
             Prediksi dan evaluasi model:
                \begin{lstlisting}[language=Python]
    y_pred_if = model_if.predict(X_test)
    y_pred_svm = model_svm.predict(X_test)
    
    from sklearn.metrics import classification_report
    print(classification_report(y_pred_if, y_pred_if))  
    print(classification_report(y_pred_svm, y_pred_svm))
                \end{lstlisting}
           
    \end{enumerate}
    
    Sistem di-{\it deploy} ke Streamlit Cloud yang mana mendukung {\it hosting} {\it framework} Streamlit secara gratis dan memungkinkan sistem dapat diakses secara {\it online} di mana pun.


% Sub-bab 2.17 node js
\section{Node.js}\label{sec:node_js}

Node.js adalah lingkungan eksekusi (\textit{runtime environment}) untuk JavaScript yang bersifat \textit{open-source} dan lintas platform, dibangun di atas mesin JavaScript V8 milik Chrome. Node.js menggunakan model I/O yang bersifat \textit{non-blocking} dan \textit{event-driven} yang menjadikannya ringan dan efisien, terutama untuk aplikasi yang berjalan secara \textit{real-time} pada perangkat terdistribusi (Tilkov \& Vinoski, 2010).

Dalam konteks penelitian ini, Node.js dipilih sebagai platform utama pengembangan sistem {data capture} dan integrasi model. Pemilihan ini didasarkan pada kebutuhan kompatibilitas dengan sistem \texttt{srusun.id} yang digunakan oleh pengelola Apartemen The Jarrdin, serta untuk memastikan integrasi yang mulus dengan modul-modul lain yang dikembangkan dalam ekosistem yang sama. Node.js memungkinkan penanganan permintaan data penyewa secara asinkronus, sehingga proses input dan validasi data penyewa dapat dilakukan dengan cepat tanpa membebani kinerja server secara keseluruhan.

Untuk mendukung fungsionalitas sistem, Node.js didukung oleh ekosistem pustaka yang luas melalui NPM (\textit{Node Package Manager}). Beberapa pustaka utama yang digunakan dalam penelitian ini meliputi \texttt{express} untuk pengelolaan server dan rute API, \texttt{mysql2/promise} untuk koneksi asinkronus ke basis data melalui mekanisme \textit{connection pool}, \texttt{jsonwebtoken} untuk keamanan autentikasi lintas platform, \texttt{multer} untuk menangani unggahan berkas (\textit{file upload}), \texttt{json2csv} untuk keperluan ekspor data laporan, serta \texttt{cors} untuk menangani akses lintas domain. Berikut langkah-langkah implementasi sistem \textit{backend} menggunakan Node.js.

\begin{enumerate}
    \item \textbf{Persiapan lingkungan dan instalasi pustaka}
    \begin{enumerate}
        \item Menginisialisasi proyek Node.js dan menginstal pustaka yang dibutuhkan untuk server dan basis data.
        \item Menginstal dependensi utama melalui terminal.
        \begin{lstlisting}[language=bash]
npm init -y
npm install express mysql2 dotenv path body-parser cors multer json2csv jsonwebtoken
        \end{lstlisting}
        \item Menyiapkan variabel lingkungan (\textit{environment variables}) pada berkas \texttt{.env} untuk menyimpan konfigurasi basis data jarak jauh (\textit{remote}) dan kunci rahasia secara aman.
    \end{enumerate}
        
    % \newpage
    \item \textbf{Konfigurasi server dan koneksi basis data}
    \begin{enumerate}
        \item Membuat konfigurasi server menggunakan \textit{framework} Express.js dan mengatur \textit{middleware} untuk keamanan dan format pertukaran data.
        \begin{lstlisting}[language=JavaScript, caption={Inisialisasi Server dan \textit{Middleware}}]
const express = require('express');
const cors = require('cors');
const bodyParser = require('body-parser');
const jwt = require('jsonwebtoken');

const app = express();

const corsOptions = {
    origin: (origin, callback) => {
        // Validasi domain yang diizinkan
        callback(null, true); 
    },
    credentials: true,
    methods: ['GET', 'POST', 'PUT', 'PATCH', 'DELETE', 'OPTIONS'],
    allowedHeaders: ['Content-Type', 'Authorization']
};

app.use(cors(corsOptions));
app.use(bodyParser.json());
app.use(bodyParser.urlencoded({ extended: true }));
        \end{lstlisting}
             
        % \newpage
        \item Membangun koneksi terpusat ke basis data MySQL (SRusun) menggunakan \textit{Connection Pool} untuk mengelola lalu lintas kueri secara efisien.
        \begin{lstlisting}[language=JavaScript, caption={Konfigurasi \textit{Connection Pool} Basis Data}]
const mysql = require('mysql2/promise');
require('dotenv').config();

const pool = mysql.createPool({
    host: process.env.SRUSUN_MYSQL_HOST,
    user: process.env.SRUSUN_MYSQL_USER,
    password: process.env.SRUSUN_MYSQL_PASSWORD,
    database: process.env.SRUSUN_MYSQL_DATABASE,
    waitForConnections: true,
    connectionLimit: 10,
    queueLimit: 0
});
        \end{lstlisting}
    \end{enumerate}
        
    \item \textbf{Implementasi API {data capture}} \\
    Membuat \textit{endpoint} API (\texttt{POST /api/rentals}) untuk menerima input data penyewa dari antarmuka pengguna (\textit{frontend}). Implementasi ini menggunakan mekanisme \texttt{\textit{database transaction}} untuk menjamin integritas data (ACID), di mana penyimpanan data profil penyewa dan pencatatan riwayat aktivitas (\textit{audit trail}) dilakukan dalam satu kesatuan proses yang saling bergantung.

    \begin{lstlisting}[language=JavaScript, caption={Implementasi Transaksi API Data Capture}]
app.post('/api/rentals', checkAuth, async (req, res) => {
    const conn = await pool.getConnection();
    try {
        const {
            nama, nik, status_pasutri, status_kewarganegaraan, 
            jenis_sewa, unit_number, metode_pembayaran, metode_lain, 
            tanggal_checkin, waktu_checkin, tanggal_checkout
        } = req.body;

        // Normalisasi NIK dan penentuan agen
        const normalizedNik = String(nik || '').replace(/\D/g, '');
        let agent_email = req.user.email;
        if (isAdmin(req.user) && req.body.agent_email) {
            agent_email = req.body.agent_email;
        }

        const lama_menginap = calculateDays(tanggal_checkin, tanggal_checkout);
        const db_status_pasutri = status_pasutri === 'Ya' ? 'Menikah' : 'Belum Menikah';

        // Memulai transaksi basis data
        await conn.beginTransaction();

        // Menyimpan data profil ke tabel rentals
        const sql = `INSERT INTO rentals (
            nama, nik, status_pasutri, status_kewarganegaraan, 
            jenis_sewa, unit_number, metode_pembayaran, metode_lain, 
            tanggal_checkin, waktu_checkin, tanggal_checkout, lama_menginap, user_email
        ) VALUES (?, ?, ?, ?, ?, ?, ?, ?, ?, ?, ?, ?, ?)`;

        const params = [
            nama, normalizedNik, db_status_pasutri, status_kewarganegaraan,
            jenis_sewa, unit_number, metode_pembayaran, 
            metode_pembayaran === 'Others' ? metode_lain : null,
            tanggal_checkin, waktu_checkin, tanggal_checkout, lama_menginap, agent_email
        ];

        const [result] = await conn.query(sql, params);
        const newRentalId = result.insertId;

        // Menyimpan rekam jejak pembuatan awal ke tabel rental_logs
        const logQuery = `INSERT INTO rental_logs (action, field_changed, new_value, email, rental_id) VALUES ?`;
        const logEntries = [['create', 'all', JSON.stringify(req.body), req.user.email, newRentalId]];
        await conn.query(logQuery, [logEntries]);

        // Menyelesaikan transaksi jika seluruh kueri berhasil
        await conn.commit();
        res.json({ success: true, message: "Data penyewa berhasil disimpan." });

    } catch (error) {
        // Membatalkan seluruh perubahan jika terjadi galat pada salah satu tahapan
        if (conn) await conn.rollback();
        console.error("Error saving rental:", error);
        res.status(500).json({ success: false, message: 'Gagal menyimpan data penyewa.' });
    } finally {
        if (conn) conn.release();
    }
});
    \end{lstlisting}
\end{enumerate}

Data yang telah tersimpan secara terstruktur dalam basis data SRusun kemudian diakses dan diolah lebih lanjut menggunakan bahasa pemrograman Python untuk tahapan klasifikasi.

